% ---------------------------------------------------------------------------
% Chapitre 11 — Fonctions de transfert, interpréteur d'expressions et stores
% Sources : notes/05-domains-transfer-stores.md (intégralement),
% src/domains/transfer/state_value.rs, src/domains/interp/{interpreter,
% callbacks,cfg}.rs, src/domains/stores/*.rs, src/domains/mod.rs (trait
% Transfer), src/engine/{fixpoint,cfg_analyzer}.rs (appelants) ;
% ADR-002, 009, 010, 012, 015, 017, 020.
% Sorties observées au commit e67b10a avec target/debug/reactant sur les
% fichiers /tmp/ch11/*.tsx (copies des exemples du dossier 05 et exemples
% propres au chapitre). Valeurs abstraites relevées avec la sonde temporaire
% du dossier 05 (/tmp/dom05/target/debug/probe05, hors dépôt : modes
% « intra » = analyze_component, « prog » = analyze_program) ; IR imprimé
% avec la sonde d'impression du chapitre 7 (~/.cache/irdump-target/debug/irdump).
% ---------------------------------------------------------------------------
\chapter{Fonctions de transfert, interpréteur d'expressions et stores}
\label{chap:transfert-stores}

\begin{objectifs}
  \item Savoir ce que calcule l'interpréteur abstrait d'un pas : la valeur d'une expression de l'IR (\code{eval_state_value}) et l'effet d'une instruction sur l'état abstrait (\code{exec_stmt}), et connaître le chemin exact des appels, du moteur jusqu'à l'évaluateur.
  \item Connaître l'environnement \code{AbstractEnv} et les quatre stores (\code{StateStore}, \code{MemoStore}, \code{SharedStateStore}, \code{Heap}) : leur clé, leur valeur par défaut ($\Top$ ou $\Bot$), leur mode de mise à jour (forte ou faible), leur portée.
  \item Savoir évaluer à la main, opération par opération, n'importe quelle expression de l'IR, et savoir ce que le lowering a déjà réécrit avant que l'évaluateur ne la voie.
  \item Comprendre comment un appel de setter, un \emph{functional updater}, un callback \code{.then} ou \code{setTimeout}, une fonction locale ou un composant enfant font bouger l'état (ADR-009, ADR-010, ADR-012).
  \item Savoir où la soundness est garantie par construction, où elle est abandonnée par décision, et reconnaître les quatorze défauts observés au commit \snapshotCommit{} (D1 à D14).
\end{objectifs}

\prerequis{\cref{chap:interpretation-abstraite} (\cref{def:treillis}, \cref{def:point-fixe}, \cref{def:widening}, \cref{def:soundness}, mises à jour fortes et faibles) ; \cref{chap:ir} (\cref{def:slot}, \cref{def:site}, \cref{def:cfg}) ; \cref{chap:lowering-cfg} (diamants, extraction des hooks, handlers) ; \cref{chap:treillis-simples} (intervalles, \code{StrConst}, opérateurs) ; \cref{chap:statevalue-stabilite} (\cref{def:statevalue}, \cref{def:stabilite}, conversion côté lecture, \code{AnalysisCtx}).}

Les trois chapitres précédents ont construit des \emph{valeurs} abstraites :
des intervalles, des ensembles de chaînes, un produit \code{StateValue} dont
chaque composante suit une sorte de JavaScript, et un treillis de stabilité
qui dit si une référence change d'un rendu à l'autre. Il manque la pièce qui
les met en mouvement : la fonction qui, devant une expression \code{n + 1} ou
une instruction \code{setN(n + 1)}, calcule la valeur abstraite ou fait
évoluer l'état. C'est l'objet de ce chapitre, la \og sémantique d'un
pas\fg{} de l'analyse. Le \cref{chap:cfg-analyzer-dominance} et le
\cref{chap:point-fixe-composant} décriront \og l'itération des pas\fg{} :
dans quel ordre on visite les blocs, quand on joint, quand on applique le \emph{widening}, quand
on s'arrête.

Le sous-système vit dans trois répertoires de \file{src/domains/}, qui
correspondent à trois questions :
\begin{itemize}
  \item \file{transfer/} : que \emph{vaut} une expression ? C'est
    l'évaluateur \code{eval_state_value}, propre au domaine
    \code{StateValue}, avec le calcul des mémos et l'inlining des
    composants enfants ;
  \item \file{interp/} : que \emph{fait} une instruction ? C'est
    l'interpréteur générique en \code{T: Transfer} : liaisons, écritures de
    champ, appels de setter, callbacks, corps de fonctions ;
  \item \file{stores/} : \emph{où} range-t-on ce que l'on sait ?
    L'environnement local, l'état React, les mémos, le tas, l'état partagé
    entre composants.
\end{itemize}
La progression suit la difficulté : on part d'un compteur, on apprend à
évaluer ses expressions, puis à exécuter ses instructions, puis à descendre
dans ses callbacks, son tas et ses enfants ; on termine par une étude de cas
des défauts de soundness observés, qui mobilise tout le reste.

% ===========================================================================
\section{Le problème : évaluer \code{n + 1} sans exécuter le programme}
\label{sec:transfert-stores-probleme}

\subsection{Deux compteurs}

Voici deux composants presque identiques.

\begin{lstlisting}[language=TypeScript,caption={Deux compteurs (\file{/tmp/ch11/counter.tsx})},label={lst:transfert-stores-counter}]
import { useState, useEffect } from "react";

export function Counter() {
  const [n, setN] = useState(0);
  useEffect(() => {
    setN(n + 1);
  });
  return <div>{n}</div>;
}

export function Settle() {
  const [m, setM] = useState(0);
  useEffect(() => {
    setM(5);
  });
  return <div>{m}</div>;
}
\end{lstlisting}

Tous deux ont un effet sans tableau de deps, donc exécuté après chaque
rendu, et cet effet appelle le setter de l'état. \code{Counter} boucle : chaque
effet écrit une valeur neuve, React refait un rendu, l'effet se relance,
indéfiniment. \code{Settle} ne boucle pas : le premier effet écrit \code{5},
le rendu suivant relance l'effet, qui écrit encore \code{5}, et React
abandonne ce second rendu parce que \code{Object.is(5, 5)} est vrai
(\cref{chap:react-modele}). L'analyseur fait bien la différence :

\begin{sortie}
$ reactant check counter.tsx --trace --show-clean
  Counter  (2 hooks)  counter.tsx
    warn   infinite-loop  [hook:0]  (line 5:2)  this effect keeps pushing state `n` to new values on every run. Potential infinite render loop
       → state `n` is written here [hook:1] (line 5:2)
       → the abstract value of state `n` kept growing and was widened at iteration 3
  Settle  (2 hooks)  counter.tsx  ✓

⚠  1 warning(s) across 1 file(s).
\end{sortie}

La règle \regle{infinite-loop} (\cref{chap:infinite-loop}) ne fait que lire
un fait établi plus tôt : \og la valeur abstraite de \code{n} a continué de
croître et a subi un widening\fg{}. Ce chapitre explique comment ce fait est
calculé. Pour cela, l'analyseur doit répondre à quatre questions, sans jamais
exécuter le programme :
\begin{enumerate}
  \item dans l'effet, que vaut \code{n} ? C'est une affaire
    d'\emph{environnement} (\cref{sec:transfert-stores-env}) ;
  \item que vaut \code{n + 1} ? C'est l'\emph{évaluateur d'expressions}
    (\cref{sec:transfert-stores-eval}) ;
  \item où va la valeur passée à \code{setN}, et comment se combine-t-elle
    avec les précédentes ? Ce sont les \emph{stores} et la \emph{mise à jour
    faible} (\cref{sec:transfert-stores-stores,sec:transfert-stores-instruction}) ;
  \item quand le corps de l'effet est-il exécuté, et que se passe-t-il s'il
    appelle un callback, une fonction locale ou un composant enfant ? Ce sont
    les \emph{points d'entrée}, la \emph{pré-passe} des callbacks, le
    \emph{tas} et l'\emph{inter-composants}
    (\cref{sec:transfert-stores-callbacks,sec:transfert-stores-tas,sec:transfert-stores-inter}).
\end{enumerate}

\subsection{Ce que voit l'interpréteur}

L'interpréteur ne voit jamais le TSX : il voit l'IR produite par le lowering
(\cref{chap:lowering-cfg}). Voici celle de \code{Counter}, imprimée par la
sonde d'impression du \cref{chap:lowering-cfg} :

\begin{sortie}
=== component Counter (param props) ===
render_cfg:
  B0:
      let n = StateVal(0)  @4:9
      let setN = StateSetter(0)  @4:12
      expr HookMarker(1, Undefined)  @5:2
      => return Native<div>({#1 })[n]
hooks:
  #0 State init=0  @4:8
  #1 Effect deps=Absent  @5:2
    B0:
        expr setN((n Add 1))  @6:4
        => return undefined
\end{sortie}

Le rendu est un bloc unique de trois instructions et un terminateur
\code{Return}. La déstructuration de \code{useState} a déjà été réécrite :
\code{n} est lié à la lecture du slot $0$ (\code{StateVal(0)}), \code{setN}
à son setter (\code{StateSetter(0)}) ; l'appel \code{useEffect} n'a laissé
qu'un marqueur de site d'appel (\code{HookMarker}). Le corps de l'effet est
un CFG à part, rangé dans la table des hooks. L'IR ne connaît que quatre
instructions, définies dans \file{src/ir/stmt.rs} : \code{Let} et
\code{Assign} (lier une variable), \code{MemberWrite} (écrire un champ,
\code{obj.f = v}) et \code{ExprStmt} (une expression évaluée pour ses
effets, \code{expr e} dans l'impression). Il n'y a pas d'instruction
\code{return} : le retour est le terminateur du bloc.

\subsection{Le calcul, à la main}

Donnons-nous, pour l'instant sans justification, trois règles :
(i) la lecture \code{StateVal(0)} rend le contenu courant du store d'état ;
(ii) \code{n + 1} se calcule sur les intervalles, $[a,b] + [1,1] = [a+1,b+1]$ ;
(iii) un appel \code{setN(v)} \emph{joint} $v$ au contenu du slot, il ne
l'écrase pas. Le store est amorcé par la valeur initiale, $[0,0]$. Le moteur
enchaîne alors des passes \og rendu puis effets\fg{} jusqu'à ce que l'état ne
bouge plus (\cref{tab:transfert-stores-iterations}).

\begin{table}[htbp]\centering\small
\begin{tabular}{lllll}
\toprule
passe & store au début & \code{n} lu & \code{n + 1} & store après la passe \\
\midrule
amorçage & --- & --- & --- & $[0,0]$ \\
1 & $[0,0]$ & $[0,0]$ & $[1,1]$ & $[0,0] \join [1,1] = [0,1]$ \\
2 & $[0,1]$ & $[0,1]$ & $[1,2]$ & $[0,2]$ \\
3 & $[0,2]$ & $[0,2]$ & $[1,3]$ & $[0,3]$, puis widening : $[0,+\infty]$ \\
4 & $[0,+\infty]$ & $[0,+\infty]$ & $[1,+\infty]$ & $[0,+\infty]$ : stable, arrêt \\
\bottomrule
\end{tabular}
\caption{Le point fixe de \code{Counter}, calculé à la main. Le widening
(\cref{def:widening}) intervient à la troisième itération, c'est-à-dire au
seuil par défaut \code{widen_threshold = 3} du moteur.}
\label{tab:transfert-stores-iterations}
\end{table}

Pour \code{Settle}, la première passe donne $[0,0] \join [5,5] = [0,5]$ et la
deuxième n'ajoute rien : le moteur s'arrête après une itération. La sonde
temporaire du dossier de notes, qui appelle directement
\code{analyze_component} et imprime les stores convergés, le confirme :

\begin{sortie}
-- Counter (intra, iterations=3)
   state[0] = number[0, inf]   widened=true
   env[n] = Val(number[0, inf])
-- Settle (intra, iterations=1)
   state[0] = number[0, 5]   widened=false
   env[m] = Val(number[0, 5])
\end{sortie}

\begin{react}[Le bail-out de \code{setState}]
  React compare la nouvelle valeur d'un état à l'ancienne par
  \code{Object.is} et n'ordonne pas de nouveau rendu si elles sont égales.
  C'est ce qui rend terminante la boucle de \code{Settle} : après le premier
  \code{setM(5)}, les suivants ne changent plus rien. Dans l'analyse, le
  point fixe abstrait est l'image de ce bail-out : une écriture qui
  n'ajoute rien au join ne relance pas la boucle.
\end{react}

Les règles (i) à (iii) sont exactement celles du code, avec des nuances que
le chapitre va détailler. Retenons dès maintenant la troisième, parce
qu'elle structure tout le reste : le store d'état contient le \emph{join de
toutes les valeurs écrites}, initialiseur compris. Il ne dit pas \og ce que
vaut l'état maintenant\fg{} mais \og ce que l'état peut valoir à un moment
quelconque\fg{}. Un appel de setter \emph{peut} avoir lieu ; il ajoute un
comportement possible, il n'en retire aucun.

\subsection{Le contrat : le trait \code{Transfer}}

Le moteur ne connaît pas \code{StateValue} : il est générique en un
paramètre \code{T: Transfer} (\srcline{src/engine/cfg_analyzer.rs}{34}). Le
trait fixe les quatre opérations dont il a besoin :

\rustsnippet[label={lst:transfert-stores-trait}]{src/domains/mod.rs}{114}{144}{le contrat entre le moteur et un domaine (la quatrième méthode, \code{recompute_memo}, suit jusqu'à la ligne 161)}

\begin{itemize}
  \item \code{eval_expr} : la valeur abstraite d'une expression. L'environnement
    est pris en lecture seule (\code{&}), l'\code{AnalysisCtx} en écriture
    (\code{&mut}) ; le \cref{chap:statevalue-stabilite} a décrit ce paquet
    (\code{component}, \code{state}, \code{memo}, \code{heap}, \code{query},
    \code{inter}) et pourquoi l'environnement en est séparé : sa mutabilité
    diffère d'une opération à l'autre.
  \item \code{exec_stmt} : l'effet d'une instruction, en place, sur
    l'environnement et les stores.
  \item \code{exec_expr_effects} : les effets d'une expression évaluée
    \og pour ses effets\fg{}, soit une instruction \code{expr;}, soit le
    corps d'une flèche concise (\code{() => setN(1)}), que le lowering range
    dans un terminateur \code{Return}. Le moteur l'appelle sur chaque
    terminateur \code{Return} qu'il rencontre
    (\src{src/engine/cfg_analyzer.rs}{84}{86}).
  \item \code{recompute_memo} : la valeur d'un \code{useMemo} ou d'un
    \code{useCallback}, recalculée après chaque passe de rendu
    (\cref{sec:transfert-stores-memo}).
\end{itemize}

Il n'existe qu'une implémentation, \code{StateValueTransfer}, une structure
vide : toute la variation vient du contexte.

\rustsnippet[label={lst:transfert-stores-svt}]{src/domains/transfer/state_value.rs}{24}{54}{l'unique implémentation du trait}

Le partage des rôles est net : l'évaluation des expressions est
\emph{propre au domaine} (\code{eval_state_value}, dans \file{transfer/}),
l'exécution des instructions est \emph{entièrement déléguée} à
l'interpréteur générique de \file{interp/}
(\code{exec_stmt_with_callbacks}, \code{exec_expr_effects}). Cet
interpréteur n'est pourtant complet que pour \code{StateValue} : le tas et le
store partagé rangent des \code{StateValue} quel que soit le domaine, et
passent par les ponts \code{as_state_value} / \code{from_state_value} du
trait \code{AbstractDomain} (\cref{chap:treillis-simples}), qui rendent
\code{None} et $\Bot$ pour tout autre domaine
(\srcline{src/domains/interp/interpreter.rs}{168}). Un second domaine perdrait
silencieusement les captures de fermetures et les écritures de champ.

\subsection{Le chemin des appels}

La \cref{fig:transfert-stores-flot} donne le chemin complet, du moteur à
l'évaluateur. Il servira de carte pour tout le chapitre.

\begin{figure}[htbp]\centering
\begin{tikzpicture}[font=\scriptsize,
    b/.style={bloc,font=\scriptsize,inner sep=3pt},
    m/.style={bloc,font=\scriptsize,inner sep=3pt,fill=white,draw=ra-gray}]
  \node[m] (acfg) at (0,0) {\code{analyze_cfg} (moteur)\\ \file{engine/cfg_analyzer.rs}};
  \node[b] (xs) at (-3.6,-1.4) {\code{Transfer::exec_stmt}\\ \code{exec_stmt_with_callbacks} (prof.~0)};
  \node[b] (full) at (-3.6,-2.9) {\code{exec_full_stmt}(prof.~$d$)};
  \node[b] (eff) at (3.2,-2.9) {\code{exec_expr_effects}(prof.~$d$)};
  \node[b] (core) at (-5.4,-4.4) {\code{exec_stmt_core}};
  \node[b] (pre) at (-1.2,-4.4) {\code{exec_callbacks_depth}\\ (pré-passe)};
  \node[b] (set) at (3.2,-4.4) {\code{exec_setter_call}};
  \node[b] (bind) at (-5.4,-5.9) {\code{bind_rhs}\\ \code{MemberWrite}};
  \node[b] (body) at (1.2,-5.9) {\code{exec_body_depth}($d+1$)\\ \code{exec_body_impl}};
  \node[b] (eval) at (-1.2,-7.4) {\code{eval_expr} = \code{eval_state_value}};
  \node[b] (binop) at (-5.6,-8.9) {\code{eval_binop}\\ \code{eval_unary}};
  \node[b] (field) at (-2.6,-8.9) {\code{eval_field_access}\\ \code{eval_index_access}};
  \node[b] (comp) at (0.8,-8.9) {\code{eval_comp_app}};
  \node[m] (child) at (4.4,-8.9) {\code{analyze_child}\\ (moteur, enfant)};
  \draw[fleche] (acfg) -- node[etiquette,left]{chaque instruction} (xs);
  \draw[fleche] (acfg) -- node[etiquette,right]{terminateur \code{Return}} (eff);
  \draw[fleche] (xs) -- (full);
  \draw[fleche] (full) -- node[etiquette,above]{\code{ExprStmt}} (eff);
  \draw[fleche] (full) -- node[etiquette,left]{\code{Let}, \code{Assign}, \code{MemberWrite}} (core);
  \draw[fleche] (full) -- (pre);
  \draw[fleche] (eff) -- (pre);
  \draw[fleche] (eff) -- (set);
  \draw[fleche] (core) -- (bind);
  \draw[fleche] (pre) -- node[etiquette,left]{callback} (body);
  \draw[fleche] (set) -- node[etiquette,right]{updater} (body);
  \draw[flechemay] (body.north west) to[out=150,in=-20] node[etiquette,right,pos=0.2]{récursion} (full.south east);
  \draw[fleche] (bind) -- (eval);
  \draw[fleche] (set) -- (eval);
  \draw[fleche] (body) -- node[etiquette,right]{\code{Return}} (eval);
  \draw[fleche] (eval) -- (binop);
  \draw[fleche] (eval) -- (field);
  \draw[fleche] (eval) -- (comp);
  \draw[fleche] (comp) -- node[etiquette,above]{point fixe} (child);
\end{tikzpicture}
\caption{Du moteur à l'évaluateur. Les boîtes blanches appartiennent au moteur
(\cref{chap:cfg-analyzer-dominance,chap:inter-composants}), les bleues à ce
chapitre. La flèche tiretée est la récursion des corps imbriqués : un corps
de callback ou de \emph{functional updater} est exécuté instruction par
instruction par le même \code{exec_full_stmt}, à la profondeur $d+1$,
plafonnée à \code{MAX_INLINE_DEPTH = 3}.}
\label{fig:transfert-stores-flot}
\end{figure}

\begin{invariant}{L'évaluation est pure vis-à-vis de l'état React}{eval-pure}
  \code{eval_state_value} ne met jamais à jour le store d'état. Les effets de
  bord (appel de setter, exécution d'un callback) sont tirés par
  l'interpréteur d'instructions, dans des positions syntaxiques précises
  (\cref{sec:transfert-stores-instruction}). Deux exceptions assumées :
  \code{havoc_setter_props}, qui joint $\Top$ dans un slot dont le setter
  part vers un enfant inanalysable, et \code{eval_comp_app}, qui lance
  l'analyse d'un enfant, laquelle écrit dans le store partagé
  (\cref{sec:transfert-stores-inter}).
\end{invariant}

Cet invariant vient de l'ADR-009 (\og la pré-passe tire les effets,
\code{eval} reste pur\fg{}). Il est commode : une même expression peut être
évaluée plusieurs fois (par la pré-passe, puis par la liaison) sans que
l'état n'enregistre deux écritures. Il a aussi un coût, que la
\cref{sec:transfert-stores-soundness} chiffrera : un setter appelé ailleurs
qu'en position d'effet n'écrit rien (défaut D2).

% ===========================================================================
\section{L'environnement abstrait}
\label{sec:transfert-stores-env}

\subsection{Un exemple : ce qu'on sait d'une variable}

L'environnement répond à la question \og que vaut la variable \code{x} en ce
point du programme ?\fg{}. Prenons un composant qui mélange les cas
courants : une prop, un appel à une fonction importée, une variable affectée
sur deux branches, une autre sur une seule.

\begin{lstlisting}[language=TypeScript,caption={Ce qu'on sait d'une variable (\file{/tmp/ch11/env.tsx})},label={lst:transfert-stores-env}]
import { useState } from "react";
import { format } from "./fmt";

export function Label({ flag }: { flag: boolean }) {
  const [n] = useState(0);
  const a = format(n);
  let b;
  if (flag) {
    b = 1;
  } else {
    b = 2;
  }
  let c;
  if (flag) {
    c = 3;
  }
  return <div>{a}{b}{c}{n}</div>;
}
\end{lstlisting}

La sonde imprime l'environnement de sortie du rendu :

\begin{sortie}
-- Label (intra, iterations=0)
   state[0] = number[0, 0]   widened=false
   env[flag] = Val(⊤)
   env[n] = Val(number[0, 0])
   env[a] = Val(⊤)
   env[b] = Val(number[1, 2])
   env[c] = Val(number[3, 3]|undefined)
\end{sortie}

Chaque ligne illustre une règle. \code{flag} est lu sur le paramètre de
props, qui n'est lié à rien quand le composant est analysé seul : il vaut
$\Top$. \code{a} est le résultat d'un appel, que l'évaluateur n'exécute
jamais : $\Top$ aussi (\cref{sec:transfert-stores-eval}). \code{b} reçoit
$[1,1]$ sur une branche, $[2,2]$ sur l'autre ; le bloc de jonction joint les
deux environnements, d'où $[1,2]$. \code{c} vaut \code{undefined} (le
\code{let c;} s'abaisse en \code{let c = undefined}) sur la branche qui ne
l'affecte pas et $[3,3]$ sur l'autre : le join garde les deux sortes, ce que
permet le produit (\cref{def:statevalue}).

\subsection{Quatre tables}

\rustsnippet[label={lst:transfert-stores-absenv}]{src/domains/stores/abstract_env.rs}{46}{60}{l'environnement abstrait}

L'environnement n'est pas une seule table mais quatre, toutes privées
(\cref{tab:transfert-stores-env}). Le nom \code{stabs} est un vestige de
l'ADR-002, quand la seule valeur abstraite était une stabilité.

\begin{table}[htbp]\centering\small
\begin{tabular}{p{27mm}p{45mm}p{30mm}p{36mm}}
\toprule
table & contenu & écrite par & lue par \\
\midrule
\code{stabs} & valeur abstraite de chaque variable & \code{extend} & \code{lookup}, \code{lookup_env_val} \\
\code{locs} & sites d'allocation que la variable peut désigner (\cref{sec:transfert-stores-tas}) & \code{extend_loc} (union) & \code{lookup_env_val} \\
\code{setter_bindings} & variable $\mapsto$ label du \code{useState} dont elle est le setter & \code{bind_setter} & \code{setter_label} \\
\code{callback_bindings} & variable $\mapsto$ label du \code{useCallback} qu'elle désigne & \code{bind_callback} & \code{callback_label} \\
\bottomrule
\end{tabular}
\caption{Les quatre tables de \code{AbstractEnv}.}
\label{tab:transfert-stores-env}
\end{table}

Les deux dernières sont des \emph{canaux latéraux} propres à React : savoir
qu'une variable est \emph{le} setter du slot $3$, ou \emph{le} callback du
hook $5$, permet d'exécuter ses appels (\cref{sec:transfert-stores-instruction,sec:transfert-stores-callbacks}).
La lecture principale est conservative :

\rustsnippet[label={lst:transfert-stores-lookup}]{src/domains/stores/abstract_env.rs}{78}{93}{la lecture d'une variable}

\code{lookup} rend $\Top$ pour toute variable absente. \code{lookup_env_val}
rend une \emph{vue} \code{EnvVal}, qui porte la valeur et, si la variable a
des sites, l'ensemble de ses sites ; cette vue sera détaillée avec le tas.
L'écriture \code{extend} écrase la valeur et \emph{préserve} les sites ;
\code{extend_loc} \emph{ajoute} un site, sans jamais en retirer
(\src{src/domains/stores/abstract_env.rs}{100}{108}).

\begin{definition}{Environnement abstrait}{environnement}
  Un \terme{environnement abstrait} $\rho$ associe à chaque variable $x$ une
  valeur $\rho(x) \in \code{StateValue}$, avec $\rho(x) = \Top$ pour toute
  variable non liée. Sa concrétisation est l'ensemble des environnements
  concrets $\sigma$ tels que $\sigma(x) \in \gam(\rho(x))$ pour tout $x$.
  L'ordre, le join et le widening sont ponctuels.
\end{definition}

\begin{invariant}{Une variable inconnue vaut $\Top$}{env-top}
  Les paramètres, les imports, les globales (\code{fetch}, \code{window}),
  toute liaison perdue se lisent $\Top$. Perdre une liaison de \emph{valeur}
  est donc sound. Les trois autres tables ne suivent pas cette règle : elles
  se lisent \og rien par défaut\fg{}. Perdre une liaison de setter, de
  callback ou de site n'est \emph{pas} sound, car ce qui n'est pas trouvé
  n'est pas exécuté.
\end{invariant}

Cette asymétrie est la racine de plusieurs défauts de la
\cref{sec:transfert-stores-soundness} (D1, D10) : partout où une liaison se
perd, la valeur retombe prudemment à $\Top$ mais l'écriture de setter qui en
dépendait disparaît.

\subsection{Le join}

\rustsnippet[label={lst:transfert-stores-envjoin}]{src/domains/stores/abstract_env.rs}{130}{173}{le join de deux environnements}

Variable par variable :
\begin{itemize}
  \item présente des deux côtés : join de \code{StateValue} ;
  \item présente d'un seul côté : $\Top$, ce qui est cohérent avec
    \og absente vaut $\Top$\fg{} ;
  \item sites : union ensembliste ;
  \item liaisons de setter et de callback : union \emph{à priorité gauche}
    (\code{entry(k).or_insert(v)}). Si une même variable désigne le setter
    de \code{a} sur un chemin et celui de \code{b} sur l'autre, le second
    est oublié (défaut D10).
\end{itemize}
Le widening \code{widen_to} (\src{src/domains/stores/abstract_env.rs}{175}{208})
suit le même schéma, avec \code{widen_to} ponctuel sur les valeurs et les
seuils de l'ADR-014 ; \code{analyze_cfg} l'emploie sur les arcs retour
(\src{src/engine/cfg_analyzer.rs}{97}{111}).

\begin{decision}[ADR-010 --- deux tables plutôt qu'une]
  La tentation était une seule table \code{Var} $\to$ \code{EnvVal}, où
  \code{EnvVal} vaut soit une valeur, soit un site. L'ADR-010 la refuse :
  \og a single map \code{EnvVal = Val | Loc} was tempting but
  \code{env.extend(var, val)} overwrote the \code{Loc} previously placed by
  \code{env.extend_loc(var, id)}\fg{}. La liaison \code{const f = () => ...}
  pose d'abord le site, puis la valeur ; avec une seule table, la seconde
  écriture efface la première et \code{f()} n'est plus exécutable. Deux
  tables indépendantes suppriment le conflit. \code{EnvVal} subsiste comme
  \emph{vue} de lecture, et comme type des champs d'un objet du tas.
\end{decision}

\begin{piege}
  \code{AbstractEnv::bottom()} est l'environnement vide. C'est bien le plus
  petit élément pour \code{leq}, mais ce n'est \emph{pas} l'élément neutre du
  join : comme une clé présente d'un seul côté donne $\Top$,
  \code{bottom().join(&e)} rend un environnement où toutes les variables de
  \code{e} valent $\Top$. En lecture, l'environnement vide se comporte comme
  \og tout $\Top$\fg{}. Les appelants évitent le piège en ne joignant que des
  environnements calculés : \code{analyze_cfg} insère le premier
  environnement reçu sans join (\srcline{src/engine/cfg_analyzer.rs}{98}),
  \code{exec_body_impl} filtre les prédécesseurs non encore visités.
  Symétriquement, \code{leq} lit une clé absente de droite comme $\Bot$ et
  non $\Top$ : l'asymétrie est sans conséquence aujourd'hui, car
  \code{AbstractEnv::leq} n'est appelé qu'en tests ; le moteur compare les
  environnements par égalité structurelle
  (\srcline{src/engine/cfg_analyzer.rs}{114}).
\end{piege}

% ===========================================================================
\section{L'évaluateur d'expressions}
\label{sec:transfert-stores-eval}

\subsection{Un banc d'essai}

Avant de lire le code, observons ce qu'il calcule sur une série
d'expressions courtes.

\begin{lstlisting}[language=TypeScript,caption={Des opérateurs (\file{/tmp/ch11/ex1.tsx})},label={lst:transfert-stores-ops}]
import { useState, useEffect } from "react";

export function Ops() {
  const [n, setN] = useState(0);
  const [label, setLabel] = useState("a");
  const masked = n & 7;
  const kind = typeof n;
  const text = `n=${n}`;
  const neg = ~5;
  const plus = +"5";
  const both = label + "b";
  const cmp = n < 3;
  const notTrue = !true;
  const notNum = !n;
  const nul = null + 1;
  useEffect(() => {
    setLabel(both);
  }, [both]);
  return <div>{masked}{kind}{text}{neg}{plus}{cmp}{notTrue}{notNum}{nul}</div>;
}
\end{lstlisting}

\begin{sortie}
-- Ops (intra, iterations=4)
   state[0] = number[0, 0]   widened=false
   state[1] = string   widened=true
   env[n] = Val(number[0, 0])
   env[label] = Val(string)
   env[masked] = Val(number[0, 7])
   env[kind] = Val(string{"number"})
   env[text] = Val(⊤)
   env[neg] = Val(number[-6, -6])
   env[plus] = Val(number[5, 5])
   env[both] = Val(string)
   env[cmp] = Val(boolean)
   env[notTrue] = Val(false)
   env[notNum] = Val(⊤)
   env[nul] = Val(number[1, 1])
\end{sortie}

Plusieurs réponses sont exactes (\code{~5} vaut $-6$, \code{+"5"} vaut $5$,
\code{null + 1} vaut $1$), d'autres sont des approximations sûres : un masque
\code{n & 7} tient dans $[0,7]$, une comparaison est \og un booléen\fg{}, et
trois valent $\Top$ alors que JavaScript saurait dire mieux (\code{text},
\code{notNum}). L'ensemble de chaînes de \code{label} grandit à chaque
passe (\code{"a"}, \code{"ab"}, \code{"abb"}, \ldots), dépasse le seuil de
\code{StrConst} et devient \code{string}, le $\Top$ de la composante
chaîne (\cref{chap:treillis-simples}). Toutes ces réponses sortent d'une
seule fonction.

\subsection{La fonction \code{eval_state_value}}

C'est une récursion structurelle sur l'arbre d'expression, un bras par
variante de \code{Expr} (\cref{chap:ir}). Elle évalue les sous-expressions
de gauche à droite, chacune une fois ; son coût est linéaire en la taille de
l'expression, sauf pour une application de composant (une analyse complète
de l'enfant) et pour la concaténation de deux ensembles de chaînes (un
produit cartésien, borné par le seuil de \code{StrConst}). Voici le début :

\rustsnippet[label={lst:transfert-stores-eval1}]{src/domains/transfer/state_value.rs}{106}{136}{l'évaluateur : littéraux, variables, hooks}

puis les marqueurs de hooks et les allocations :

\rustsnippet[label={lst:transfert-stores-eval2}]{src/domains/transfer/state_value.rs}{137}{163}{l'évaluateur : marqueurs, allocations, éléments}

La suite (\src{src/domains/transfer/state_value.rs}{165}{190}) évalue les
opérandes d'un \code{BinOp} ou d'un \code{UnaryOp} puis délègue à
\code{eval_binop} ou \code{eval_unary}, traite les appels, les \code{new},
les accès de membres et d'index, et traverse les annotations. La
\cref{tab:transfert-stores-eval} récapitule tous les bras ; on la détaille
ensuite par familles.

\begin{table}[htbp]\centering\small
\begin{tabular}{p{42mm}p{42mm}p{62mm}}
\toprule
expression de l'IR & valeur abstraite & remarque \\
\midrule
\code{Lit(Int(n))}, \code{Lit(Float(f))} & \code{number([n,n])} & intervalle ponctuel \\
\code{Lit(Bool(b))} & \code{boolean(True|False)} & \\
\code{Lit(String(s))} & \code{str_singleton(s)} & ensemble $\{s\}$ \\
\code{Lit(Null)}, \code{Lit(Unit)} & \code{null()}, \code{undefined()} & \code{undefined} s'abaisse en \code{Lit(Unit)} \\
\code{Var(v)} & \code{env.lookup(v)} & $\Top$ si non liée \\
\code{StateVal(l)} & \code{ctx.state.get(l)}, référence convertie & \cref{sec:transfert-stores-lecture} \\
\code{StateSetter(l)} & \code{component_setter}\newline\code{(ctx.component, l)} & le setter porte son propriétaire \\
\code{MemoVal(l)}, \code{CallbackVal(l)} & \code{ctx.memo.get(l)} & $\Top$ tant que non calculé \\
\code{HookMarker(_, Undefined)} & \code{undefined()} & hook React sans valeur suivie \\
\code{HookMarker(_, Unknown)} & $\Top$ & hook custom ni inliné ni résumé \\
\code{HookMarker(_, StableRef)} & \code{ref(Stable)} & \code{useRef} \\
\code{HookMarker(_, Summary(s))}, \code{SummaryVal(s)} & \code{summary_value(s)} & table unique des résumés \\
\code{ObjectLit}, \code{ArrayLit}, \code{FnLit}, \code{New} & \code{ref(PerRender)} & fait must : neuf à chaque rendu \\
\code{NativeElem}, \code{CompApp} & \code{ref(Stable)} & voir D8 \\
\code{BinOp}, \code{UnaryOp} & \code{eval_binop}, \code{eval_unary} & opérandes d'abord \\
\code{Call} & \code{ref(PerRender)} ou $\Top$ & allocateurs connus seulement \\
\code{FieldAccess}, \code{IndexAccess} & lecture par le tas & \cref{sec:transfert-stores-tas} \\
\code{TSAnnotated(e)} & valeur de \code{e} & l'annotation est ignorée \\
\bottomrule
\end{tabular}
\caption{Les bras de \code{eval_state_value}, \src{src/domains/transfer/state_value.rs}{106}{190}.}
\label{tab:transfert-stores-eval}
\end{table}

\subsection{Littéraux, variables et hooks}

Les littéraux sont exacts. Un entier du source devient \code{Int} quand sa
valeur est entière et petite, \code{Float} sinon ; les deux donnent un
intervalle ponctuel. Une variable se lit dans l'environnement, $\Top$ si elle
n'y est pas.

Trois bras lisent des hooks. \code{StateVal(l)} lit le store d'état ; il
applique la \emph{conversion côté lecture} de l'ADR-017, que le
\cref{chap:statevalue-stabilite} a démontrée et que la
\cref{sec:transfert-stores-lecture} replacera dans l'interpréteur.
\code{StateSetter(l)} rend un setter \emph{porteur de son propriétaire},
\code{SetterVal::One(composant, l)}, même en analyse intra-composant : c'est
ce qui permet de suivre un setter passé en prop jusqu'à l'enfant qui
l'appelle (\cref{sec:transfert-stores-inter}), et c'est aussi l'origine du
défaut D3. \code{MemoVal} et \code{CallbackVal} lisent le même
\code{MemoStore} : un \code{useCallback} n'est qu'un mémo dont la valeur est
une fonction.

Le marqueur \code{HookMarker} ancre le site d'appel d'un hook dont la
valeur n'est pas suivie. Ses quatre formes ont chacune une histoire,
racontée par les commentaires du code : un hook React sans valeur
(\code{useEffect}) rend vraiment \code{undefined} ; un hook custom non
résolu rendait autrefois \code{undefined} lui aussi, ce qui le faisait
passer pour \emph{prouvablement stable} et taisait toutes les règles de
stabilité (un faux négatif corrigé par \code{b7bc459}) ; il rend désormais
$\Top$. \code{useRef} rend le même conteneur à chaque rendu : une référence,
et une référence \code{Stable}. Un hook de bibliothèque résumé se lit comme
son résumé, par une table unique, \code{summary_value}
(\src{src/domains/transfer/state_value.rs}{192}{222}), \og so the marker and
the standalone \code{SummaryVal} can never disagree\fg{} : \code{Top} et
\code{Held} donnent $\Top$, \code{StableRef} donne \code{ref(Stable)},
\code{UnstableRef} \code{ref(PerRender)}, un \code{Wrapper} l'une ou l'autre
selon son drapeau. Deux lignes méritent l'attention. Un résumé \code{Shape}
(\code{useForm()}, dont seuls certains membres sont promis stables) vaut
$\Top$ : la promesse porte sur les membres, qui seront lus par le tas
(\cref{sec:transfert-stores-tas}). Un \code{Navigator} non stable vaut $\Top$
et non \code{PerRender} : on ne prouve pas sa fraîcheur.

\subsection{Allocations et éléments}

Un littéral d'objet, de tableau ou de fonction évalué pendant un rendu est
une référence neuve à chaque rendu : \code{ref(PerRender)} est un fait
\emph{must} (\cref{def:stabilite}). C'est le fait que la règle
\regle{always-unstable-deps} exploite. \code{new X()} alloue lui aussi
(ajout du commit \snapshotCommit{}, \issue{158}) ; ses membres restent
$\Top$. L'évaluation d'un littéral \emph{ne peuple pas} le tas : seule la
liaison d'une variable (\code{bind_rhs}) ou le passage en prop le fait.

Les éléments JSX, natifs (\code{<span/>}) ou composants (\code{<Child/>}),
valent \code{ref(Stable)}. C'est juste pour dire \og la valeur d'un élément
ne porte pas d'information utile\fg{}, mais c'est faux au sens
d'\code{Object.is} : en JavaScript, chaque rendu alloue un nouvel objet
élément. Le \cref{chap:statevalue-stabilite} a relevé ce faux négatif ; on
le retrouvera sous le nom D8.

\subsection{Appels}

Un appel n'est \emph{jamais} exécuté par l'évaluateur : sa valeur est $\Top$,
sauf pour une liste blanche d'appels qui allouent toujours.

\rustsnippet[label={lst:transfert-stores-fresh}]{src/domains/transfer/state_value.rs}{238}{259}{les appels dont le résultat est toujours neuf}

La liste est volontairement courte. Elle exclut \code{slice} et
\code{concat}, qui rendent une primitive sur une chaîne
(\code{id.slice(0, 8)}) : prétendre \code{PerRender} serait une fausse
preuve (\issue{22}, décision enregistrée). Elle exclut aussi les méthodes en
place (\code{sort}, \code{reverse}, \code{fill}, \code{Object.assign}), qui
rendent le receveur, soit le fait inverse. Les allocateurs statiques
(\code{from}, \code{keys}, \code{parse}\ldots) ne comptent que sur leur
receveur connu (\code{Array}, \code{Object}, \code{JSON}).

Deux conséquences sont à garder en tête. D'abord, la \emph{valeur de retour}
d'une fonction locale appelée n'est jamais calculée : même quand
l'interpréteur exécute son corps pour ses effets
(\cref{sec:transfert-stores-callbacks}), il jette le retour. Ensuite, les
\emph{effets} d'un appel (un setter, un callback) ne passent pas par ici :
ils sont tirés par l'interpréteur d'instructions.

\subsection{Opérateurs}

Le \cref{chap:treillis-simples} a décrit en détail \code{eval_binop} et
\code{eval_bitwise} (\src{src/domains/transfer/state_value.rs}{752}{929}) ;
on n'en garde ici que ce qu'il faut pour calculer à la main.
\begin{itemize}
  \item \code{+}, \code{-}, \code{*} : arithmétique d'intervalles quand les
    deux opérandes sont dans $\{$nombre, \code{null}$\}$, avec
    $\code{ToNumber(null)} = 0$ (vue \code{as_arith},
    \cref{chap:statevalue-stabilite}) ; pour \code{+}, concaténation exacte
    si les deux opérandes sont des ensembles finis de chaînes, \og une
    chaîne\fg{} si ce sont des chaînes mal connues ; $\Top$ dans tous les
    autres cas, y compris \code{"n=" + n}.
  \item \code{/} : toujours $\Top$ (un intervalle ne contient ni \code{NaN}
    ni $\pm$\code{Infinity}) ; \texttt{\%} et \code{**} : exacts quand
    \code{NaN} est exclu, $\Top$ sinon.
  \item comparaisons, \code{in}, \code{instanceof} : \code{boolean(Top)}.
    Rien n'est décidé, mais le résultat \emph{n'est qu'un booléen}, ce qui
    garde les composantes nombre et chaîne à $\Bot$.
  \item \code{&&}, \code{||} : $\Top$ ; en pratique ces variantes
    n'arrivent jamais ici (\cref{sec:transfert-stores-lowering}).
  \item bit à bit : toujours un nombre de la plage int32 (uint32 pour
    \code{>>>}), raffiné pour un masque constant (\code{x & 7} dans
    $[0,7]$) ou un décalage constant. Un opérande $\Bot$ rend $\Bot$ : un
    chemin tué par un raffinement doit rester mort.
\end{itemize}

Les opérateurs unaires sont propres à ce fichier :

\rustsnippet[label={lst:transfert-stores-unary}]{src/domains/transfer/state_value.rs}{941}{977}{les opérateurs unaires (début)}

\begin{itemize}
  \item \code{-x} : négation d'intervalle, $\Top$ hors $\{$nombre,
    \code{null}$\}$ ;
  \item \code{!x} : exact sur un booléen pur, $\Top$ sinon ;
  \item \code{typeof x} : une chaîne, exacte (\code{"number"},
    \code{"object"} pour \code{null}, \code{"function"} pour un setter\ldots)
    quand une seule sorte est habitée, \code{string} sinon. Cette exactitude
    fait de \code{typeof x === "string"} une garde raffinable ;
  \item \code{~x} : $-(\code{ToInt32}(x) + 1)$, bornes échangées dans int32,
    plage int32 sinon ;
  \item \code{+x} : \code{ToNumber}, pas l'identité ; un booléen donne $0$,
    $1$ ou $[0,1]$, un ensemble de chaînes donne l'enveloppe de leurs
    valeurs si \emph{toutes} se lisent comme des nombres finis
    (\code{coerce_to_number}), $\Top$ sinon ;
  \item \code{UnaryOp::Unknown} : $\Top$.
\end{itemize}
La fin de la fonction (\code{~}, \code{+} et le cas inconnu) occupe
\src{src/domains/transfer/state_value.rs}{978}{1004}, et
\code{coerce_to_number} \src{src/domains/transfer/state_value.rs}{1006}{1040}.

\begin{piege}
  Deux imprécisions surprennent à la lecture de la sonde. \code{!n} sur un
  nombre vaut $\Top$ et non \code{boolean} : la politique \og c'est toujours
  un booléen\fg{} des comparaisons n'a pas été appliquée à la négation. Un
  gabarit \code{`n=${n}`} vaut $\Top$ et non \code{string} : le lowering l'a
  réécrit en \code{"n=" + n}, et \code{+} entre une chaîne et un nombre
  tombe dans le cas \og tout le reste\fg{}. Les deux sont sound ; ils sont
  simplement plus lâches que JavaScript. La lecture des chaînes par
  \code{coerce_to_number} utilise l'analyseur de Rust
  (\code{s.trim().parse()}), plus strict que \code{ToNumber} : \code{+""}
  vaut $0$ en JavaScript mais $\Top$ ici. Là encore, c'est sound.
\end{piege}

\subsection{Membres et annotations}

\code{o.f} et \code{xs[0]} se lisent dans le tas abstrait ; on y reviendra
en \cref{sec:transfert-stores-tas}, une fois le tas présenté. Un index non
constant (\code{xs[i]}) vaut $\Top$. Une annotation de type
(\code{TSAnnotated}) est transparente : l'ADR-020 a explicitement refusé de
brancher les types TypeScript dans le domaine.

% ===========================================================================
\section{Ce que le lowering a déjà fait}
\label{sec:transfert-stores-lowering}

L'évaluateur n'a aucun bras pour les gabarits, les ternaires, \code{&&},
\code{||}, \code{??}, le chaînage optionnel, les spreads, \code{++},
\code{void}, \code{delete}, \code{await} ou les séquences. Ces constructions
n'arrivent jamais jusqu'à lui : le lowering
(\file{src/lowering/expr_lower.rs}, \cref{chap:lowering-cfg}) les a réécrites
en constructions déjà traitées. Pour prédire une valeur abstraite, il faut
donc connaître ces réécritures (\cref{tab:transfert-stores-lowering}).

\begin{table}[htbp]\centering\small
\begin{tabular}{p{30mm}p{52mm}p{64mm}}
\toprule
source & IR produite & conséquence pour la valeur \\
\midrule
\code{`q${e}r`} & \code{("q" + e) + "r"} & exact si \code{e} est une chaîne connue ; $\Top$ si c'est un nombre \\
\code{tag`...${e}...`} & \code{Call(tag, [e, ...])} & $\Top$ (appel), les lectures de \code{e} survivent \\
\code{a ? b : c} & diamant : \code{Branch(a)}, \code{let __tN = b} / \code{let __tN = c}, jonction qui lit \code{__tN} & join des deux branches \emph{par le CFG}, avec raffinement sur \code{a} \\
\code{a && b}, \code{a || b}, \code{a ?? b} & \code{let __tN = a} ; \code{Branch(__tN)} ; \code{__tN := b} dans le bloc droit & idem ; \code{b} est en position d'\code{Assign} (compte pour D2) \\
\code{a?.b}, \code{f?.(x)} & comme \code{a.b}, \code{f(x)} & la nullité du receveur n'est pas modélisée (compte pour D5) \\
\code{{ ...o, k: v }} & membre sous une clé synthétique \code{"...N"} & seuls les membres après le dernier spread sont suivis \\
\code{[...xs]} & \code{xs} gardé comme élément, arité \code{AtLeast(n)} & aucune influence sur la valeur \\
\code{f(...args)} & \code{args} émis comme instruction & un setter passé par spread n'est pas un argument (\issue{76}) \\
\code{{ [k]: v }}, \code{get x()} & clés synthétiques \code{"[computed]N"}, \code{"[accessor]N"} & membre introuvable par son nom \\
\code{i++}, \code{--i} & \code{i := i + 1} (ou \code{- 1}), valeur \code{Var(i)} & post-écriture pour les deux formes \\
\code{o.f++}, \code{delete o.f} & \code{MemberWrite} opaque / de \code{undefined} & mise à jour faible du champ \\
\code{void e}, \code{(a, b, c)} & \code{ExprStmt} pour les effets, puis la valeur & effets conservés \\
\code{await e} & \code{let __tN = e}, puis coupure de bloc & la suite est dans un bloc successeur \\
\code{this}, \code{super} & \code{SummaryVal(Top)} & $\Top$ \\
classe & \code{ObjectLit} de ses membres & référence fraîche \\
\bottomrule
\end{tabular}
\caption{Les réécritures du lowering qui déterminent la valeur calculée
(d'après \file{src/lowering/expr_lower.rs}).}
\label{tab:transfert-stores-lowering}
\end{table}

\begin{exemple}{Un ternaire n'est pas évalué par l'évaluateur}{transfert-stores-ternaire}
  L'effet \code{setB((x) => (x < 10 ? x + 1 : x))} devient, dans l'IR, un
  corps de fonction à quatre blocs :
\begin{sortie}
        expr setB(fn#8(x) {
            B0:
                => branch (x Lt 10) ? B1 : B2
            B1:
                let __t0 = (x Add 1)  @22:26
                => jump B3
            B2:
                let __t0 = x  @22:34
                => jump B3
            B3:
                => return __t0
            edges: 0->1[T] 0->2[F] 1->3[U] 2->3[U]
        })  @22:4
\end{sortie}
  L'évaluateur ne voit que \code{x Lt 10}, \code{x Add 1}, \code{x} et
  \code{__t0}. Le \og ternaire\fg{} est le join des environnements des blocs
  \code{B1} et \code{B2} au bloc \code{B3}. C'est à l'interpréteur de corps
  (\cref{sec:transfert-stores-callbacks}) de le calculer correctement ; on
  verra qu'il ne le fait pas (défaut D1).
\end{exemple}

\begin{decision}[ADR-020 n\textsuperscript{o}~1 --- garder le diamant de \code{&&}/\code{||}]
  On pourrait ajouter à l'IR un nœud plat \code{LogicalOp} évalué par
  l'évaluateur, comme \code{BinOp}. L'ADR-020 le refuse : le diamant modélise
  les \emph{effets} d'un court-circuit (dans \code{flag && setA(1)}, l'appel
  n'a lieu que sur une branche) et permet au moteur de raffiner
  l'environnement sur la condition. Le coût accepté est une valeur moins
  précise quand l'expression est lue, et la dispersion du \code{rhs} dans un
  \code{Assign}, qui a une conséquence imprévue : défaut D2. Le
  \cref{chap:lowering-cfg} note en outre que la nature des deux arêtes d'un
  court-circuit encode l'opérateur, pas la valeur de vérité ; le point fixe
  n'en dépend pas, car il raffine d'après \code{then_}/\code{else_}.
\end{decision}

% ===========================================================================
\section{Les stores}
\label{sec:transfert-stores-stores}

L'environnement décrit un \emph{point} du programme : il change d'un bloc à
l'autre, et le moteur en garde un par bloc. Ce que React conserve
\emph{d'un rendu à l'autre} (l'état, les mémos, les objets qui survivent)
ne peut pas vivre là : il faut des structures dont la portée est le
composant, voire le programme. Ce sont les \emph{stores}.

\subsection{\code{StateStore} : l'état React, sujet du point fixe}

\rustsnippet[label={lst:transfert-stores-statestore}]{src/domains/stores/state_store.rs}{7}{45}{le store d'état : lecture, mise à jour faible, fusion}

Le store associe à chaque label de \code{useState} (ou de \code{useReducer})
une \code{StateValue}. Trois propriétés le définissent.

\begin{itemize}
  \item \textbf{$\Bot$ par défaut.} Un label jamais écrit se lit $\Bot$,
    \og impossible\fg{}, à l'inverse de l'environnement. C'est correct parce
    que tout slot est \emph{amorcé} avant la boucle par la valeur de son
    initialiseur (ci-dessous) : un $\Bot$ ne peut provenir que d'un hook mal
    extrait.
  \item \textbf{Mise à jour faible.} \code{update} ne remplace jamais : elle
    joint. Un appel de setter ajoute un comportement possible (sémantique
    may). Le store contient donc, à tout instant, le join de toutes les
    valeurs écrites, initialiseur compris : c'est la \emph{vue événement} de
    l'ADR-017.
  \item \textbf{Fusion sur l'union des clés.} \code{join}, \code{widen} et
    \code{widen_to} ne diffèrent que par le combinateur passé à
    \code{merge_with}.
\end{itemize}

L'amorçage se fait dans le moteur, une fois par analyse de composant
(\src{src/engine/fixpoint.rs}{314}{353}) ; pour chaque \code{HookEntry::State},
la valeur initiale est calculée ainsi, puis jointe au slot par
\code{state.update(*label, init_val)} :

\rustsnippet[label={lst:transfert-stores-seed}]{src/engine/fixpoint.rs}{336}{348}{l'amorçage d'un slot par son initialiseur}

Deux cas. Un initialiseur paresseux \code{useState(() => e)} est exécuté
comme un corps de fonction (\code{exec_body}, \cref{sec:transfert-stores-callbacks})
et c'est son \emph{retour} qui amorce le slot : React exécute la fonction une
fois au montage et range son résultat, jamais la fonction elle-même. Tout
autre initialiseur est simplement évalué. Dans les deux cas, le contexte est
\code{AnalysisCtx::null} avec un store vide, des mémos vides et un tas neuf :
l'initialiseur voit l'environnement d'entrée du composant (constantes de
module, props liées par un parent), mais ni l'état ni le tas du composant.

La boucle du point fixe (\cref{chap:point-fixe-composant}) n'emploie le store
qu'à travers quelques méthodes : \code{bottom} pour l'état initial et les
accumulateurs de passe, \code{join} pour combiner rendu, effets, handlers et
tranche partagée, \code{leq} pour le test d'arrêt \og le nouvel état est-il
$\lleq$ l'ancien ?\fg{}, \code{widen_to} à partir de
\code{widen_threshold} itérations, \code{widen} (sans seuils) au garde-fou
des cent itérations, et \code{changed_labels} pour savoir quels slots ont été
soumis au widening (\src{src/engine/fixpoint.rs}{486}{530}).

\subsection{\code{MemoStore} : des valeurs dérivées, sans point fixe propre}

\rustsnippet[label={lst:transfert-stores-memostore}]{src/domains/stores/memo_store.rs}{27}{35}{le store des mémos : lecture et écriture}

Le \code{MemoStore} range la valeur de chaque \code{useMemo} et de chaque
\code{useCallback}. Il est l'opposé du store d'état sur les deux axes :
\begin{itemize}
  \item \textbf{$\Top$ par défaut} : un mémo pas encore calculé peut valoir
    n'importe quoi ;
  \item \textbf{mise à jour forte} : \code{set} écrase. Le mémo est recalculé
    en entier après chaque passe de rendu, à partir de l'environnement de
    sortie du rendu (\src{src/engine/fixpoint.rs}{382}{413}). Ce n'est pas un
    sujet du point fixe : la convergence du store d'état entraîne celle des
    mémos, y compris d'une chaîne mémo $\to$ mémo, d'une itération à la
    suivante.
\end{itemize}
Pendant une passe, \code{analyze_cfg} travaille sur une copie jetée à chaque
bloc (\og Memo is fixed for this pass; local mutations are discarded\fg{},
\src{src/engine/cfg_analyzer.rs}{65}{66}).

\subsection{\code{SharedStateStore} : l'état vu par tout le programme}

\rustsnippet[label={lst:transfert-stores-shared}]{src/domains/stores/shared_state_store.rs}{10}{17}{le store inter-composants (\code{get} et \code{update} suivent, l.~24--34)}

Quand un enfant appelle un setter reçu en prop, l'écriture concerne un slot
d'un \emph{autre} composant. Elle ne peut pas aller dans le store d'état de
l'enfant ; elle va dans ce store-ci, indexé par le couple
(\code{ComponentId} du propriétaire, label). Il est non générique (fixé à
\code{StateValue}), $\Bot$ par défaut, mis à jour par join, exactement comme
un \code{StateStore}. Il en existe \emph{un seul par programme}, créé par
\code{analyze_program} (\srcline{src/engine/fixpoint.rs}{714}) et partagé par
toutes les analyses, racines et enfants inlinés à toute profondeur, à toutes
les itérations. Il n'est jamais remis à zéro : c'est un accumulateur monotone
global. Chaque composant en importe sa part à chaque test de convergence :

\rustsnippet[label={lst:transfert-stores-slice}]{src/domains/stores/shared_state_store.rs}{57}{66}{la tranche d'un composant}

C'est tout le mécanisme de remontée de l'état : \og No separate
program-level fixpoint layer\fg{} (ADR-012). Pour le point fixe du parent,
une écriture faite par un enfant est indiscernable d'une écriture locale.

\subsection{\code{Heap} : le tas abstrait}

Le quatrième store range les objets et les fermetures créés par le
composant, indexés par leur site d'allocation (\cref{def:site}). Il mérite
une section à lui seul (\cref{sec:transfert-stores-tas}) ; retenons pour
l'instant qu'il a une entrée par \code{ExprId}, qu'il est mis à jour par
écrasement (sauf les champs d'objets, joints), et qu'il n'entre pas dans le
test de convergence.

\subsection{Synthèse}

\begin{definition}{Store}{store}
  Un \terme{store} est une table persistante d'une passe à l'autre, dont la
  clé désigne une entité du programme React plutôt qu'une variable : un slot
  d'état (\code{StateStore}), un mémo (\code{MemoStore}), un slot d'un
  composant quelconque (\code{SharedStateStore}) ou un site d'allocation
  (\code{Heap}). Un store est caractérisé par sa clé, sa valeur par défaut
  ($\Bot$ si toute clé légitime est amorcée, $\Top$ sinon), son mode de mise
  à jour (forte, par écrasement, ou faible, par join) et sa portée (un
  composant ou le programme). L'environnement \code{AbstractEnv}, indexé par
  variable et attaché à un point du programme, n'est pas un store.
\end{definition}

\begin{table}[htbp]\centering\small
\begin{tabular}{p{24mm}p{21mm}p{20mm}p{24mm}p{16mm}p{23mm}}
\toprule
& clé & défaut & mise à jour & point fixe ? & portée \\
\midrule
\code{AbstractEnv} & \code{Var} & $\Top$ (valeur) ; rien (setter, callback, site) & \code{extend} forte ; \code{extend_loc} union & par bloc & un point de programme \\
\code{StateStore} & \code{HookLabel} & $\Bot$ & \code{update} = join & \textbf{oui} & un composant \\
\code{MemoStore} & \code{HookLabel} & $\Top$ & \code{set} forte & non, recalculé & un composant \\
\code{SharedStateStore} & \code{(ComponentId,}\newline\code{HookLabel)} & $\Bot$ & \code{update} = join & via \code{slice} & le programme \\
\code{Heap} & \code{ExprId} & absent & \code{insert} forte ; champs joints & non & un composant (copie aux enfants) \\
\bottomrule
\end{tabular}
\caption{L'environnement et les quatre stores (\cref{def:store}).}
\label{tab:transfert-stores-stores}
\end{table}

\begin{figure}[htbp]\centering
\begin{tikzpicture}[font=\scriptsize,
    s/.style={bloc,font=\scriptsize,text width=36mm,minimum height=15mm}]
  \node[s] (env) at (0,1.2) {\textbf{\code{AbstractEnv}}\\ clé \code{Var} ; défaut $\Top$\\ \code{extend} (forte)\\ \emph{un par bloc}};
  \node[s] (state) at (-4.6,-2.4) {\textbf{\code{StateStore}}\\ clé label ; défaut $\Bot$\\ \code{update} = $\join$ (faible)\\ \emph{sujet du point fixe}};
  \node[s] (memo) at (0,-2.4) {\textbf{\code{MemoStore}}\\ clé label ; défaut $\Top$\\ \code{set} (forte)\\ \emph{recalculé après le rendu}};
  \node[s] (heap) at (4.6,-2.4) {\textbf{\code{Heap}}\\ clé \code{ExprId} ; absent\\ \code{insert} (forte), champs $\join$\\ \emph{insensible au flot}};
  \node[s] (shared) at (-4.6,-5.6) {\textbf{\code{SharedStateStore}}\\ clé (composant, label) ; $\Bot$\\ \code{update} = $\join$ (faible)\\ \emph{un par programme}};
  \node[draw=ra-gray,dashed,rounded corners,fit=(state)(memo)(heap),inner sep=5pt] (compo) {};
  \node[etiquette,anchor=north east] at (compo.south east) {portée : un composant};
  \draw[fleche] (state.north) |- node[etiquette,below,pos=0.75]{lecture \code{StateVal(l)}} (env.west);
  \draw[fleche] (heap.north) |- node[etiquette,above,pos=0.75]{lecture \code{o.f}, sites} (env.east);
  \draw[fleche] ([xshift=-5mm]memo.north) -- node[etiquette,left]{\code{MemoVal(l)}} ([xshift=-5mm]env.south);
  \draw[fleche] ([xshift=5mm]env.south) -- node[etiquette,right]{recalcul} ([xshift=5mm]memo.north);
  \draw[fleche] ([xshift=-12mm]env.south) -- node[etiquette,left,pos=0.35]{appel de setter} ([xshift=12mm]state.north);
  \draw[flechemay] ([yshift=4mm]env.west) -- node[etiquette,above,pos=0.6]{setter d'un autre composant} ++(-5.3,0) |- (shared.west);
  \draw[fleche] (shared.north) -- node[etiquette,right]{\code{slice}, à chaque test de convergence} (state.south);
  \node[etiquette,align=left,right=4mm of shared] {l'enfant y écrit les slots de ses ancêtres ;\\ chaque parent en importe sa tranche};
\end{tikzpicture}
\caption{La carte des stores. Les flèches vers l'environnement sont des
lectures faites par l'évaluateur ; la flèche \og appel de setter\fg{} est
l'écriture faible d'un setter local (\cref{sec:transfert-stores-instruction}),
la flèche tiretée celle d'un setter reçu d'un ancêtre ; la flèche
\og recalcul\fg{} est la mise à jour des mémos depuis l'environnement de
sortie du rendu. Le store partagé est le seul dont la portée dépasse le
composant.}
\label{fig:transfert-stores-carte}
\end{figure}

\begin{decision}[ADR-002 --- trois stores séparés, et ce qu'il en reste]
  L'ADR-002 a posé le principe : \og A unified store would force all hooks
  into the same fixpoint\fg{}. L'état est seul sujet du point fixe, parce
  que seule sa mise à jour déclenche un nouveau rendu ; les mémos sont
  dérivés de leurs deps ; les refs sont triviales. Le principe tient, mais le
  détail a changé sans que l'ADR soit marquée remplacée : le domaine est
  devenu le produit \code{StateValue} (ADR-015) ; le \code{MemoStore} range
  une valeur et non un couple (deps, valeur), et il est recalculé à
  \emph{chaque} itération, pas une fois après stabilisation ; le
  \code{RefStore} n'a jamais existé, \code{useRef} est un marqueur
  \code{StableRef} ; et deux stores sont apparus, le tas (ADR-010) et le
  store partagé (ADR-012). L'alternative refusée, un store unifié, aurait
  obligé les mémos et les refs à converger comme l'état.
\end{decision}

\begin{piege}
  Les valeurs par défaut sont inversées d'un store à l'autre :
  l'environnement et les mémos rendent $\Top$, l'état et le store partagé
  rendent $\Bot$. La cohérence tient à l'amorçage. Une règle ou un contributeur
  qui lirait un label d'état jamais amorcé obtiendrait $\Bot$, c'est-à-dire
  un chemin mort, et en tirerait des conclusions fausses. Autre confusion
  classique : \code{state_store.get(l)} (la vue événement, le join des
  écritures) n'est pas ce que lit un rendu ; le rendu lit
  \code{StateVal(l)}, converti (\cref{sec:transfert-stores-lecture}).
\end{piege}

% ===========================================================================
\section{L'instruction : liaisons, écritures de champ, setters}
\label{sec:transfert-stores-instruction}

\subsection{Le squelette : pré-passe, puis cœur}

\code{Transfer::exec_stmt} délègue à \code{exec_stmt_with_callbacks}, qui
appelle \code{exec_full_stmt} à la profondeur $0$
(\src{src/domains/interp/interpreter.rs}{29}{36}). Tout se décide là :

\rustsnippet[label={lst:transfert-stores-fullstmt}]{src/domains/interp/interpreter.rs}{69}{99}{une instruction : pré-passe des callbacks, puis cœur}

Une instruction s'exécute toujours en deux temps. D'abord la
\emph{pré-passe} \code{exec_callbacks_depth} parcourt les expressions de
l'instruction et exécute les callbacks \og dans le cycle\fg{} qu'elle y
trouve (\code{.then}, \code{setTimeout}, \code{map}\ldots,
\cref{sec:transfert-stores-callbacks}). Ensuite le \emph{cœur} fait ce que
l'instruction dit : lier, écrire un champ, ou tirer les effets d'une
expression. Une instruction \code{expr;} passe entièrement par
\code{exec_expr_effects}, qui contient sa propre pré-passe.

\subsection{Lier une variable : \code{bind_rhs}}

\code{Let} et \code{Assign} partagent une seule séquence. Le commentaire du
code rappelle pourquoi : \og the arms had drifted, dropping the Assign
aliases (FN)\fg{} (\src{src/domains/interp/interpreter.rs}{143}{149}) ;
\code{let s = setX} et \code{s = setX} doivent se comporter à l'identique.

\rustsnippet[label={lst:transfert-stores-bind1}]{src/domains/interp/interpreter.rs}{214}{245}{lier une variable (1/2) : setters, callbacks, allocations}

\rustsnippet[label={lst:transfert-stores-bind2}]{src/domains/interp/interpreter.rs}{246}{292}{lier une variable (2/2) : résumés, alias, chaînes de membres, valeur}

Pas à pas, pour \code{var = rhs} :
\begin{enumerate}
  \item si \code{rhs} est un setter (\code{StateSetter(l)}), \code{var} est
    lié au label \code{l} dans la table des setters ; si c'est un
    \code{useCallback} (\code{CallbackVal(l)}), dans la table des callbacks ;
  \item si \code{rhs} est un littéral de fonction, \code{var} reçoit son
    site, et la fermeture est allouée dans le tas avec ses captures
    (\code{alloc_fn}, \cref{sec:transfert-stores-tas}) ;
  \item si \code{rhs} est un littéral d'objet, la carte de ses membres est
    rangée dans le tas sous son site, et \code{var} reçoit ce site ;
  \item si \code{rhs} est un résumé \code{Shape} de bibliothèque
    (\code{const f = useForm()}), même chose avec les membres promis ;
  \item si \code{rhs} est une autre variable (\code{let y = x}), \code{y}
    hérite des sites, de la liaison de setter et de la liaison de callback de
    \code{x} : c'est ainsi qu'un alias de setter reste un setter ;
  \item si \code{rhs} est une chaîne de membres (\code{let f = props.onClick}),
    \code{var} reçoit les sites que \code{resolve_locs} trouve au bout de la
    chaîne ;
  \item enfin, la valeur est évaluée et rangée (\code{env.extend}).
\end{enumerate}

L'ordre a une conséquence subtile. Pour un littéral de fonction, le site est
posé \emph{avant} l'allocation, et \code{alloc_fn} capture l'environnement
\emph{avant} que \code{var} ne soit liée par valeur : une fonction récursive
\code{const f = () => f()} capture \code{f} avec sa valeur antérieure, soit
$\Top$ à la première liaison. Autre conséquence : un tableau littéral n'est
jamais rangé dans le tas, et un objet littéral n'y est rangé que s'il est
\emph{directement} lié ; un objet imbriqué (\code{{ inner: { f } }}) n'a
pas d'entrée propre (défaut D6).

\subsection{Écrire un champ : une mise à jour faible}

\rustsnippet[label={lst:transfert-stores-memberwrite}]{src/domains/interp/interpreter.rs}{155}{200}{l'écriture de champ, mise à jour faible du tas}

\code{obj.f = v} ne change pas l'identité de \code{obj} : c'est ce qui en
fait une mutation. L'écriture peut n'avoir lieu que sur un chemin, et
\code{obj} peut désigner plusieurs sites : on \emph{joint}, on ne remplace
jamais. Les valeurs sont toujours jointes ; les sites ne survivent que si
l'ancienne et la nouvelle valeur en ont (une valeur simple qui écrase une
fermeture ne laisse rien à poursuivre). Trois limites : seul
\code{Var.champ = v} est traité (pas \code{a.b.c = v}, pas
\code{o[k] = v}) ; un champ \emph{nouveau} (cas \code{(None, new)}) reçoit
la seule nouvelle valeur, sans \code{undefined} ; et le tas n'étant pas
sensible au flot, l'écriture faite sur une branche est vue sur l'autre. Les
deux dernières se combinent dans le défaut D9.

\subsection{Tirer les effets d'une expression}

\rustsnippet[label={lst:transfert-stores-effects}]{src/domains/interp/interpreter.rs}{114}{126}{la définition unique de \og exécuter une expression pour ses effets\fg{}}

Trois effets, dans cet ordre : la pré-passe des callbacks, l'appel de setter
éventuel, et l'inlining d'un composant enfant si l'expression en est un
(\code{return <Child/>}). C'est la définition unique de \og exécuter une
expression pour ses effets\fg{}, partagée par \code{ExprStmt} et par le
terminateur \code{Return} du moteur. Le cœur est l'appel de setter :

\rustsnippet[label={lst:transfert-stores-setter}]{src/domains/interp/interpreter.rs}{341}{392}{l'appel de setter : slot local, puis slot d'un propriétaire}

\paragraph{Premier bras : un setter lié localement.} Si le callee est une
variable liée à un setter (\code{env.setter_label(name)}), on calcule la
valeur écrite et on la joint au slot (\code{ctx.state.update}). La valeur
écrite dépend de l'argument :
\begin{itemize}
  \item un littéral de fonction, le \emph{functional updater}
    \code{setX(c => ...)} : son premier paramètre est lié à la valeur
    courante du store (\code{ctx.state.get(label)}), son corps est exécuté
    par \code{exec_body_depth} à la profondeur $d+1$, et sa \emph{valeur de
    retour} est la valeur écrite ;
  \item tout autre argument : sa valeur ;
  \item aucun argument : $\Top$ (JavaScript écrirait \code{undefined} ;
    $\Top$ est sound).
\end{itemize}

\paragraph{Second bras : un setter porteur de propriétaire.} Si la valeur du
callee est un setter \code{SetterVal::One(c, l)} et qu'une analyse
inter-composants est en cours (\code{ctx.inter} présent), on joint la valeur
du premier argument à \code{shared_state[(c, l)]}. C'est ainsi que
\code{props.onChange(42)}, dans un enfant, écrit l'état du parent. Ce bras
est \emph{toujours} évalué, que le premier ait tiré ou non. Comme
\code{StateSetter(l)} s'évalue en \code{SetterVal::One(composant courant, l)}
(\cref{sec:transfert-stores-eval}), un setter local le tire aussi : pour un
argument ordinaire, l'écriture dans le store partagé est redondante mais
inoffensive ; pour un \emph{functional updater}, elle ne l'est pas (défaut
D3).

\begin{react}[\emph{Functional updater} et file d'attente]
  \code{setX(prev => f(prev))} reçoit la valeur la plus récente, après les
  mises à jour déjà en file, alors que \code{setX(x + 1)} lit la valeur
  capturée par la fermeture au moment du rendu. Plusieurs appels d'un même
  handler sont regroupés (\emph{batching}). L'analyse lie le paramètre de
  l'updater au join du store, qui sur-approxime toute valeur en file ; elle
  ne cherche pas à modéliser l'ordre des mises à jour.
\end{react}

\begin{exemple}{Dérouler un \emph{functional updater}}{transfert-stores-updater}
  Le test unitaire \code{functional_updater_increments_state}
  (\src{src/domains/transfer/state_value.rs}{1480}{1521}) construit à la main
  l'IR de \code{setN(c => c + 1)} sur un store où le slot $0$ vaut $[5,5]$.
  Le paramètre \code{c} est lié à $[5,5]$ ; le corps rend $[6,6]$ ; le slot
  devient $[5,5] \join [6,6] = [5,6]$. L'ancienne valeur reste : la mise à
  jour est faible. Sur un vrai programme,
  \code{useEffect(() => { setA((x) => x + 1); })} mène le slot au widening comme le
  compteur de la \cref{sec:transfert-stores-probleme} ; la sonde, en mode
  intra-composant, donne \code{state[0] = number[0, inf] widened=true} pour
  \code{UpdaterNoDeps} (\file{/tmp/ch11/ex7.tsx}). On verra en
  \cref{sec:transfert-stores-soundness} que le même programme, analysé par
  le CLI, ne donne pas ce résultat.
\end{exemple}

\begin{invariant}{Où un setter tire}{setter-positions}
  Un appel de setter écrit dans le store d'état en trois positions
  exactement : une instruction \code{expr;} (\code{ExprStmt}), un
  terminateur \code{Return} de CFG analysé par le moteur, et le terminateur
  \code{Return} d'un corps exécuté par \code{exec_body_impl}. Nulle part
  ailleurs : un appel dans le membre droit d'une liaison, en argument d'un
  autre appel, ou dans le bras droit d'un \code{&&} (qui devient un
  \code{Assign}) n'écrit rien.
\end{invariant}

C'est une asymétrie avec le reste du moteur. La relation
\relation{slot_writers} du \cref{chap:relations} voit un setter appelé dans
\emph{n'importe quelle} position d'expression depuis \issue{130} ; le store
d'état, lui, ne le voit qu'aux trois positions de l'invariant. Le défaut D2
en découle.

\begin{remarque}
  \code{useReducer} passe par le même chemin : le lowering extrait
  \code{useReducer(reducer, init)} comme un état ordinaire, en sautant le
  réducteur, et \code{dispatch} est lié comme un setter. \code{dispatch(a)}
  écrit donc l'\emph{action} \code{a} dans le slot, pas
  \code{reducer(s, a)} (défaut D12).
\end{remarque}

% ===========================================================================
\section{La conversion côté lecture}
\label{sec:transfert-stores-lecture}

Le store range la vue \emph{événement} : le join des valeurs écrites. Un
rendu, lui, lit la vue \emph{inter-rendus} : entre deux rendus, la valeur
d'un état ne peut changer qu'aux appels de son setter. Le bras
\code{StateVal(l)} de l'évaluateur passe de l'une à l'autre
(\src{src/domains/transfer/state_value.rs}{122}{134}) : si la composante
\code{reference} de la valeur stockée est non-$\Bot$, elle est
\emph{remplacée} par \code{Versioned({(composant, l)})} ; les autres
composantes passent intactes. Le \cref{chap:statevalue-stabilite} a
démontré la soundness de cette conversion sous l'hypothèse \og aucun setter
n'est appelé pendant le rendu\fg{}, dont la violation a son propre
diagnostic (\regle{setter-in-render}). Ce qui compte ici est sa
\emph{position} : c'est le seul point de conversion du projet (\og The
conversion happens in \textbf{one place}\fg{}, ADR-017), et tout ce qui est
en aval en hérite gratuitement : liaisons, deps, props, lectures de champs.

\begin{exemple}{Les deux vues sur un même slot}{transfert-stores-deux-vues}
  Dans \file{/tmp/ch11/ex17.tsx}, un état objet
  \code{const [obj, setObj] = useState({ a: 1 })} est réécrit par un effet
  \code{setObj({ a: 2 })}. La sonde imprime le store et la variable :
\begin{sortie}
   state[1] = ref(PerRender)   widened=false
   env[obj] = Val(ref(Versioned({(ComponentId(4294967295), 1)})))
\end{sortie}
  Le store dit vrai : on n'y a écrit que des littéraux frais. La variable
  dit vrai aussi : d'un rendu à l'autre, \code{obj} ne change qu'aux appels
  de \code{setObj}. Le composant est ici \code{ComponentId::SYNTHETIC}
  ($2^{32}-1$), l'identité d'un composant analysé seul.
\end{exemple}

La lecture d'un champ prolonge la conversion : si le receveur est versionné
et que le tas ne connaît pas le champ, \code{eval_field_access} garde les
étiquettes de version au lieu de rendre un $\Top$ complet
(\cref{sec:transfert-stores-tas}). Un champ d'un état objet ne change, lui
aussi, qu'aux écritures de cet état, sous l'hypothèse qu'on ne mute pas
l'état en place (la règle \regle{state-mutation} traite la violation).

% ===========================================================================
\section{Les callbacks : classe de déclenchement, pré-passe, corps}
\label{sec:transfert-stores-callbacks}

\subsection{Le piège de l'\code{onClick}}

Jusqu'ici, un effet n'écrivait l'état que par un appel direct. Le code réel
passe par des callbacks :

\begin{lstlisting}[language=TypeScript,caption={Cinq façons d'appeler un setter depuis un effet (\file{/tmp/ch11/ex9.tsx})},label={lst:transfert-stores-callbacks}]
import { useState, useEffect } from "react";
import { myHelper } from "./lib";

export function Callbacks() {
  const [user, setUser] = useState(null);
  const [n, setN] = useState(0);
  const [clicks, setClicks] = useState(0);
  const [h, setH] = useState(0);
  const [d, setD] = useState(0);
  useEffect(() => {
    fetch("/api").then((u) => setUser(u));
    const tick = () => setN(42);
    setTimeout(tick, 100);
    window.addEventListener("click", () => setClicks(7));
    myHelper(() => setH(9));
    const f1 = () => setD(1);
    const f2 = () => f1();
    const f3 = () => f2();
    const f4 = () => f3();
    f4();
  }, []);
  return <div>{user}{n}{clicks}{h}{d}</div>;
}
\end{lstlisting}

Si l'interpréteur ignore les corps de fonctions, les cinq écritures sont
invisibles : c'était la situation avant l'ADR-009, où seul le
\emph{functional updater} était descendu. Si au contraire il descend
\emph{uniformément} dans tous les callbacks, il exécute le handler de
\code{click} comme s'il tournait à chaque effet, et un
\code{addEventListener("click", () => setCount(c => c + 1))} parfaitement
correct devient une \og boucle infinie\fg{}.

\begin{react}[Quand un callback s'exécute-t-il ?]
  \code{arr.map(cb)} exécute \code{cb} tout de suite. \code{.then(cb)},
  \code{setTimeout(cb)}, \code{queueMicrotask(cb)},
  \code{requestAnimationFrame(cb)} l'exécutent plus tard, mais \emph{en
  conséquence} de l'effet qui les planifie : ils font partie du cycle
  automatique rendu, commit, effets. Un \code{addEventListener} ou un
  \code{onClick} ne l'exécute que sur un événement externe, hors de ce
  cycle.
\end{react}

\subsection{La classe de déclenchement}

La réponse de l'ADR-009 est de classer chaque appel par \emph{ce qui
déclenche} ses arguments fonctions :

\rustsnippet[label={lst:transfert-stores-trigger}]{src/domains/interp/callbacks.rs}{6}{56}{la classe de déclenchement d'un appel}

\begin{itemize}
  \item \code{Setter} : le callee est un setter lié. Son argument fonction
    est un \emph{functional updater}, exécuté par \code{exec_setter_call} ;
    la pré-passe ne doit pas le descendre, sinon il s'exécuterait deux fois.
  \item \code{InCycle} : fonctions d'ordre supérieur synchrones (\code{map},
    \code{forEach}, \code{reduce}, \code{Array.from}\ldots) et planifications
    asynchrones (\code{then}, \code{catch}, \code{finally}, les minuteurs).
    Leurs arguments fonctions sont exécutés.
  \item \code{Subscription} : \code{addEventListener},
    \code{removeEventListener}. Non descendus ici ; le lowering relève
    \code{addEventListener(nom, fn)} dans un effet comme un
    \code{HookEntry::Handler}, que le moteur analyse comme un \emph{point
    d'entrée} séparé, exclu du signal de widening.
  \item \code{Unknown} : tout le reste, en particulier les helpers et hooks
    custom. Non descendus.
\end{itemize}

La reconnaissance est \emph{nominale} : c'est le nom du callee qui décide,
sans aucune information de type. Un \code{x.map(cb)} dont \code{x} n'est pas
un tableau est tout de même \og dans le cycle\fg{} : on exécute un corps qui
ne l'aurait peut-être pas été, ce qui peut produire un faux positif, jamais
un faux négatif.

\begin{decision}[ADR-009 --- \code{Unknown} n'est pas descendu]
  Descendre un callee inconnu serait le choix sound : rien ne prouve que
  \code{myHelper(cb)} n'appelle pas \code{cb} dans le cycle. L'ADR-009 le
  refuse : \og for a linter, false positives are more costly than false
  negatives\fg{}, car descendre produirait un FP sur chaque enveloppe de
  souscription custom (\code{useInterval}, \code{useEventCallback}).
  \og We accept the FN\fg{} (\issue{19}) ; la table de politique est un
  bouton réglable. C'est l'un des rares cas documentés où la précision
  l'emporte sur l'invariant \og faux négatifs interdits\fg{}
  (\cref{def:soundness}), et il n'est déclaré par aucune Info. Écarts avec
  l'ADR : elle distinguait \code{InCycleSync} et \code{InCycleDeferred},
  que le code fusionne (même politique), et le code ajoute \code{Setter}.
\end{decision}

\subsection{La pré-passe}

\rustsnippet[label={lst:transfert-stores-prepass}]{src/domains/interp/interpreter.rs}{477}{540}{la pré-passe des callbacks}

Pas à pas :
\begin{enumerate}
  \item \textbf{Garde de profondeur.} À la profondeur
    \code{MAX_INLINE_DEPTH = 3}, on s'arrête. En analyse inter-composants,
    on l'enregistre dans \code{callback_depth_capped}, ce qui donnera une
    Info \regle{analysis-limit} et suspendra les vérifications du composant ;
    en analyse intra, le plafond est silencieux.
  \item \textbf{Appel.} On classe le callee, on descend le receveur (pour
    les chaînes \code{a.then(x).then(y)}), puis chaque argument :
    \begin{itemize}
      \item un littéral de fonction, si l'appel est \code{InCycle} : son corps
        est exécuté par \code{exec_body_depth} à la profondeur $d+1$, dans
        l'environnement \emph{vivant} du site d'appel (capture naturelle d'une
        fermeture en ligne), paramètres à $\Top$ (la valeur résolue d'une
        promesse, l'élément d'un tableau sont inconnus) ;
      \item une variable, si l'appel est \code{InCycle} : c'est le cas
        \emph{B5} de l'ADR-010 (\code{setTimeout(tick)}), résolu par le tas
        (\code{exec_var_callback}, \cref{sec:transfert-stores-tas}) ;
      \item un littéral de fonction d'une autre classe : rien ;
      \item tout autre argument : descente récursive (un appel imbriqué est
        traité à son tour).
    \end{itemize}
  \item \textbf{Appel direct d'une fonction locale.} Si le callee est
    \code{Unknown} \emph{et} une variable, on tente de l'exécuter par le tas :
    c'est le cas \emph{B6} (\code{load()}). Une fonction importée n'a pas de
    site et n'est pas exécutée.
  \item \textbf{Composant enfant.} Pour un \code{CompApp} rencontré en
    profondeur (un enfant JSX dans un autre élément), pré-passe sur ses
    props puis évaluation, ce qui déclenche l'inlining
    (\cref{sec:transfert-stores-inter}).
  \item \textbf{Tout le reste} : descente générique dans les
    sous-expressions, qui ne descend jamais dans un littéral de fonction.
\end{enumerate}

\begin{invariant}{La pré-passe n'entre jamais dans un corps}{prepass-corps}
  La pré-passe ne descend pas \emph{dans} un littéral de fonction ; un corps
  ne s'exécute que par \code{exec_body_depth}. Sinon, comme
  \code{exec_body_impl} rappelle \code{exec_full_stmt}, qui rappelle la
  pré-passe, chaque corps serait exécuté deux fois.
\end{invariant}

\subsection{Exécuter un corps : \code{exec_body_impl}}

\rustsnippet[label={lst:transfert-stores-body}]{src/domains/interp/interpreter.rs}{394}{469}{l'exécution d'un corps de fonction imbriqué}

Tout corps imbriqué passe par là : callbacks, \emph{functional updaters},
initialiseurs paresseux, arguments de hooks custom, fonctions locales
inlinées. L'algorithme est volontairement simple :
\begin{enumerate}
  \item ordonner les blocs par un tri topologique, post-ordre inversé d'un
    parcours en profondeur depuis l'entrée (\code{topo_sort},
    \src{src/domains/interp/cfg.rs}{5}{11}) ; les arcs retour sont ignorés,
    l'en-tête d'une boucle passe avant son corps, les blocs inatteignables ne
    sont jamais visités ;
  \item pour chaque bloc, dans cet ordre, une seule fois : calculer son
    environnement comme le join de ceux de ses prédécesseurs, exécuter ses
    instructions par \code{exec_full_stmt} (à la même profondeur), puis
    traiter le terminateur ;
  \item un \code{Return} tire les effets de son expression, comme une
    instruction, puis joint sa valeur à la valeur de retour ; un saut ou un
    branchement propage l'environnement vers les successeurs ;
  \item s'il existe un arc retour, la valeur de retour est jointe à $\Top$.
\end{enumerate}

C'est le \emph{second interpréteur de CFG} du projet. Le premier,
\code{analyze_cfg} (\cref{chap:cfg-analyzer-dominance}), sert au rendu, aux
effets et aux handlers : une worklist, avec raffinement sur les gardes et
widening. Celui-ci n'a ni l'un ni l'autre : une seule passe, les deux
branches d'un \code{if (true)} jointes, le corps d'une boucle exécuté une
fois avec l'environnement d'avant la boucle. Les setters d'une boucle tirent
donc bien, mais la valeur portée par la boucle n'est vue qu'à sa première
itération (\issue{21}, limite déclarée dans \file{docs/limitations.md}).

\begin{piege}
  Deux interpréteurs de force inégale (\issue{12}) : une même expression
  peut être évaluée avec une précision très différente selon qu'elle est
  écrite dans un effet ou dans le \code{setTimeout} de cet effet. Et le
  second a un défaut de propagation (D1) : l'étape 2 joint les
  environnements \emph{d'entrée} des prédécesseurs, pas leurs sorties. On
  le démontrera en \cref{sec:transfert-stores-soundness}.
\end{piege}

\subsection{Retour sur l'exemple}

La sonde, en mode programme (le chemin du CLI), donne :

\begin{sortie}
-- ComponentId(0) (iterations=1)
   state[0] = ⊤   widened=false   shared=⊤
   state[1] = number[0, 42]   widened=false   shared=number[42, 42]
   state[2] = number[0, 7]   widened=false   shared=number[7, 7]
   state[3] = number[0, 0]   widened=false   shared=⊥
   state[4] = number[0, 0]   widened=false   shared=⊥
stats: cache_hits=0 cache_misses=0 depth_capped={ComponentId(0)} unknown_refs={}
\end{sortie}

\begin{table}[htbp]\centering\small
\begin{tabular}{p{44mm}p{70mm}p{26mm}}
\toprule
appel & mécanisme & slot \\
\midrule
\code{fetch().then(u => setUser(u))} & \code{then} est \code{InCycle}, le littéral est exécuté, \code{u} vaut $\Top$ & \code{user} = $\Top$ \\
\code{setTimeout(tick, 100)} & \code{InCycle}, argument variable : B5, le site de \code{tick} est dans le tas & \code{n} = $[0,42]$ \\
\code{addEventListener("click", ...)} & \code{Subscription}, non descendu ; relevé en handler et analysé comme point d'entrée & \code{clicks} = $[0,7]$ \\
\code{myHelper(() => setH(9))} & callee importé : \code{Unknown}, pas de site, pas de B6 & \code{h} = $[0,0]$ (FN accepté) \\
\code{f4()} $\to$ \code{f3} $\to$ \code{f2} $\to$ \code{f1} & B6 à chaque niveau ; profondeurs 1, 2, 3 : plafond & \code{d} = $[0,0]$ \\
\bottomrule
\end{tabular}
\caption{Les cinq appels de \cref{lst:transfert-stores-callbacks}.}
\label{tab:transfert-stores-callbacks}
\end{table}

La colonne \code{shared} montre que les écritures locales sont aussi versées
au store partagé en mode programme (le second bras de
\code{exec_setter_call}) ; c'est inoffensif ici, aucun argument n'étant un
updater. Le CLI déclare le plafond et suspend les vérifications ; il ne dit
rien de \code{myHelper} :

\begin{sortie}
$ reactant check ex9.tsx --info
  Callbacks  (7 hooks)  ex9.tsx
    info   analysis-limit  callback inlining reached the depth cap (3), so deeper HOF chains are not descended (FN possible on nested callbacks)
    warn   missing-cleanup  [hook:5]  (line 14:4)  this effect calls `window.addEventListener` but returns no cleanup. The registration is repeated every time the effect re-runs (and on every mount, twice under StrictMode) and nothing ever undoes it; return a function that tears it down
    suspended  analysis-limit  10 passing check(s) withheld: the analysis was truncated in this component, so they are not guaranteed
\end{sortie}

% ===========================================================================
\section{Le tas par site d'allocation}
\label{sec:transfert-stores-tas}

\subsection{Le besoin}

L'environnement dit ce que \emph{vaut} une variable. Deux questions lui
échappent : \og quel corps exécuter quand on appelle \code{tick} ?\fg{} (B5,
B6) et \og que vaut le membre \code{onClear} de l'objet \code{bag} ?\fg{}.

\begin{lstlisting}[language=TypeScript,caption={Membres, sites et fermetures (\file{/tmp/ch11/ex10.tsx}, extrait du corps de \code{Heapy})},label={lst:transfert-stores-heapy}]
  const [q, setQ] = useState("");
  const onClear = useCallback(() => setQ(""), []);
  const bag = { onClear, label: "x", n: 1 };
  const fromBag = bag.onClear;
  const { label } = bag;
  const fresh = new Map();
  const items = [1, 2].map((x) => x * 2);
  const nested = { inner: { f: () => setQ("n") } };
  const g = nested.inner.f;
  useEffect(() => {
    g();
  }, [fromBag, items]);
\end{lstlisting}

\code{bag} est un objet neuf à chaque rendu, mais son membre \code{onClear}
est \emph{la même} fonction à chaque rendu (un \code{useCallback} aux deps
vides). Si \code{bag.onClear} héritait de la stabilité du conteneur, la dep
\code{fromBag} serait signalée à tort ; c'était le faux positif de
\issue{88}. Il faut donc une mémoire \emph{par membre}.

\subsection{Sites, valeurs du tas, vue \code{EnvVal}}

\rustsnippet[label={lst:transfert-stores-heap}]{src/domains/stores/heap.rs}{17}{34}{le tas abstrait et ses valeurs}

La clé est le \emph{site d'allocation} (\cref{def:site}) : l'\code{ExprId}
que le lowering attribue à chaque \code{ObjectLit}, \code{ArrayLit},
\code{FnLit} et \code{New}, stable d'une itération à l'autre et unique dans
un composant. Les sites synthétiques (l'objet de props d'un enfant) viennent
de \code{ExprId::fresh()}, un compteur qui démarre à $10^9$. Une valeur du
tas est soit une fermeture (\code{Fn} : paramètres, CFG du corps partagé par
\code{Arc}, valeurs capturées de ses variables libres), soit un objet
(\code{Obj} : une carte membre $\mapsto$ \code{EnvVal}).

\code{EnvVal} (\src{src/domains/stores/abstract_env.rs}{10}{27}) est la vue
d'une liaison : \code{Val(v)}, une valeur seule, ou \code{Loc { ids, val }},
un ensemble de sites \emph{et} une valeur. Le commentaire du code insiste :
un site et une valeur ne sont pas des alternatives. Rendre $\Top$ pour la
valeur d'un \code{Loc}, comme autrefois, faisait lire une prop fonction comme
\og inconnue\fg{} dans l'enfant, alors que le parent l'avait prouvée
\code{PerRender}.

\subsection{Allouer une fermeture : un point unique}

\rustsnippet[label={lst:transfert-stores-allocfn}]{src/domains/stores/heap.rs}{45}{70}{l'allocation d'une fermeture et de ses captures}

\code{alloc_fn} est appelé par les trois bras qui peuvent lier un littéral de
fonction (\code{Let}/\code{Assign}, membre d'objet, \code{MemberWrite}) et par
l'évaluation des props d'un enfant. Il calcule les variables libres du corps
et capture leur \emph{valeur} dans l'environnement courant ; une variable
libre non liée est capturée comme $\Top$. Ni les sites ni les liaisons de
setter des variables capturées ne sont capturés ; un setter reste retrouvable
parce que sa \emph{valeur} est un \code{SetterVal::One(composant, label)}, que
le second bras de \code{exec_setter_call} sait appeler.

\subsection{Suivre une chaîne de membres : un walker unique}

\rustsnippet[label={lst:transfert-stores-resolve}]{src/domains/stores/heap.rs}{134}{165}{les sites que désigne une expression}

\code{resolve_locs} ne connaît que trois formes : une variable (ses sites),
une annotation (transparente), et un accès de champ (les sites du membre, sur
chaque objet que le receveur peut désigner). Un index, un appel, un littéral
imbriqué rendent \code{None}. C'est le seul walker du projet : l'évaluateur
de champs, la liaison \code{let} et la carte des membres d'un objet
l'appellent tous, et le commentaire rappelle que les anciens cas particuliers
\code{Expr::Var} \og each stopped at the first segment\fg{}.

\subsection{Lire un membre}

\rustsnippet[label={lst:transfert-stores-field}]{src/domains/transfer/state_value.rs}{598}{628}{la lecture d'un membre par le tas}

On résout les sites du receveur ; sur chaque site qui est un objet
\emph{possédant} le champ, on prend la valeur du membre ; on joint. Si aucun
site ne répond, on retombe sur $\Top$, en gardant les étiquettes de version
du receveur (\cref{sec:transfert-stores-lecture}). Un index entier constant
se lit comme le membre de même nom (\code{xs[0]} comme le membre
\code{"0"}, \src{src/domains/transfer/state_value.rs}{641}{651}) : c'est ce
qui permet à un hook résumé qui rend un tuple
(\code{const [value, setValue] = useAtom(a)}) d'exposer ses positions.

\begin{piege}
  L'étape \og on joint les objets qui ont le champ\fg{} ignore ceux qui ne
  l'ont pas, au lieu d'y joindre \code{undefined} ou $\Top$. Si le receveur
  peut désigner un objet sans ce membre (une autre branche, un objet issu
  d'un spread, un receveur nullable lu par \code{?.}), la réponse est
  sous-approximée. La doc de \code{obj_members} affirme le contraire
  (\og an absent member falls back to the container's own value --- sound,
  just imprecise\fg{}) ; ce n'est vrai que si \emph{aucun} site n'a le
  membre. C'est le défaut D5.
\end{piege}

\begin{figure}[htbp]\centering
\begin{tikzpicture}[font=\scriptsize,
    e/.style={cfgbloc,text width=45mm},
    h/.style={cfgbloc,text width=58mm}]
  \node[etiquette] at (-2.4,0.9) {environnement de sortie du rendu};
  \node[etiquette] at (4.6,0.9) {tas du composant};
  \node[e] (ebag) at (-2.4,0) {bag $\mapsto$ Loc \{2\}, ref(PerRender)};
  \node[e] (efrom) at (-2.4,-0.8) {fromBag $\mapsto$ Val ref(Stable)};
  \node[e] (elabel) at (-2.4,-1.6) {label $\mapsto$ Val string\{"x"\}};
  \node[e] (enested) at (-2.4,-2.6) {nested $\mapsto$ Loc \{6\}, ref(PerRender)};
  \node[e] (eg) at (-2.4,-3.4) {g $\mapsto$ Val $\Top$};
  \node[h] (h2) at (4.6,-0.8) {\#2 $\mapsto$ Obj \{\\ \quad onClear $\mapsto$ Val ref(Stable)\\ \quad label $\mapsto$ Val string\{"x"\}\\ \quad n $\mapsto$ Val number[1, 1] \}};
  \node[h] (h6) at (4.6,-3.0) {\#6 $\mapsto$ Obj \{\\ \quad inner $\mapsto$ Val ref(PerRender) \}\\ (le littéral intérieur n'a pas d'entrée)};
  \draw[flechemust] (ebag.east) -- (h2.west |- ebag.east);
  \draw[flechemust] (enested.east) -- (h6.west |- enested.east);
  \draw[flechemay] (h2.west |- efrom.east) -- node[etiquette,above]{\code{bag.onClear}} (efrom.east);
  \draw[flechemay] (h2.west |- elabel.east) -- node[etiquette,below]{\code{__obj.label}} (elabel.east);
  \draw[flechemay] (h6.west |- eg.east) -- node[etiquette,above]{\code{.inner} : pas de site} (eg.east);
\end{tikzpicture}
\caption{Le tas de \code{Heapy}, d'après la sonde. Les flèches pleines vont
d'une variable à ses sites ; les tiretées sont les lectures de membres qui
ont produit une liaison. \code{bag.onClear} lit la valeur \emph{propre} du
membre, \code{Stable}, et non celle du conteneur ; \code{nested.inner.f}
s'arrête faute de site pour le littéral intérieur, et \code{g} vaut
$\Top$.}
\label{fig:transfert-stores-tas}
\end{figure}

La sonde confirme la \cref{fig:transfert-stores-tas} :

\begin{sortie}
-- Heapy (intra, iterations=0)
   state[0] = string{""}   widened=false
   env[bag] = Loc { ids: {ExprId(2)}, val: ref(PerRender) }
   env[fromBag] = Val(ref(Stable))
   env[label] = Val(string{"x"})
   env[fresh] = Val(ref(PerRender))
   env[items] = Val(ref(PerRender))
   env[nested] = Loc { ids: {ExprId(6)}, val: ref(PerRender) }
   env[g] = Val(⊤)
\end{sortie}

et le CLI en tire deux avertissements justes, et un silence faux :

\begin{sortie}
$ reactant check ex10.tsx
  Heapy  (4 hooks)  ex10.tsx
    warn   always-unstable-deps  [hook:2]  (line 13:2)  this effect depends on `items`, a new reference every render, so `Object.is` always differs and the effect re-runs on every render regardless of the other deps
       (1 trace step(s), rerun with --trace)
    warn   missing-deps  [hook:2]  var:g  (line 13:2)  `g` is used in this effect but not in its deps array, and its value may change between renders
       (1 trace step(s), rerun with --trace)
\end{sortie}

\code{fromBag} n'est pas signalé (sa stabilité est prouvée), \code{items} l'est
(sa fraîcheur est certaine), \code{g} l'est parce qu'il vaut $\Top$. Mais
\code{g()} n'a pas été exécuté dans l'effet : \code{setQ("n")} est manqué,
\code{q} reste \code{{""}}, et le mémo \code{[q]} du composant est jugé
stable alors qu'il changera (défaut D6).

\subsection{Exécuter une fonction par son nom : B5 et B6}

\rustsnippet[label={lst:transfert-stores-varcb}]{src/domains/interp/interpreter.rs}{544}{583}{exécuter un callback désigné par une variable}

C'est ici que le tas rejoint la pré-passe. Pour une variable liée à des
sites, on exécute \emph{tous} les corps possibles (\og All bodies are
executed and their effects joined --- correct by over-approximation\fg{},
ADR-010), dans l'environnement vivant, complété par les captures (qui
l'emportent) et les paramètres à $\Top$. Pour une variable liée à un
\code{useCallback}, le corps est demandé à \code{QueryContext::callback_body}
(il vit dans la table des hooks, pas dans le tas). Pour une variable sans
site ni callback (un import, un paramètre), on ne fait rien.

Trois programmes isolent la frontière (\file{/tmp/ch11/ex11.tsx}) :
\code{Flat} lie \code{g = o.f} avec \code{o = { f: () => setQ("n") }} ;
\code{Nested} lie \code{g = o.inner.f} avec un objet imbriqué ;
\code{MethodCall} appelle directement \code{o.f()}. Seul le premier voit
l'écriture :

\begin{sortie}
-- Flat (intra, iterations=1)
   state[0] = string{"", "n"}   widened=false
-- Nested (intra, iterations=0)
   state[0] = string{""}   widened=false
-- MethodCall (intra, iterations=0)
   state[0] = string{""}   widened=false
\end{sortie}

Dans \code{Flat}, \code{resolve_locs(o.f)} trouve le site du littéral, B6
exécute le corps. Dans \code{Nested}, le littéral intérieur n'a pas de site.
Dans \code{MethodCall}, le callee est un \code{FieldAccess} : B6 ne traite que
les callees \code{Var}.

\begin{decision}[ADR-010 --- un tas par site d'allocation]
  Pour relier un nom à un corps, l'ADR-010 choisit l'abstraction classique
  par site d'allocation : un \code{ExprId} par nœud qui alloue, une table
  \code{ExprId} $\to$ valeur du tas, un environnement à deux tables, et
  \emph{un seul tas} partagé par toutes les passes et itérations d'un
  composant (sans quoi le \code{tick} alloué dans le rendu serait inconnu de
  l'effet). Limites déclarées : un callee \code{Unknown} sans site est un FN ;
  une profondeur au-delà de trois est un FN signalé ; les membres avant un
  spread et les accesseurs manquent à la carte. Écarts avec le code : les
  fermetures ont gagné leurs captures (ADR-012), les objets une carte
  d'\code{EnvVal} (\issue{88}), et la variante \code{HeapValue::Arr},
  \og reserved --- future array domain\fg{}, n'a jamais été écrite.
\end{decision}

\begin{piege}
  Le tas n'est pas un treillis qu'on itère. \code{Heap::join},
  \code{widen} et \code{leq} existent
  (\src{src/domains/stores/heap.rs}{80}{131}) mais ne sont appelés nulle part ;
  le moteur manipule un seul tas mutable par composant, passé à tous les
  blocs, toutes les passes, toutes les itérations. Trois conséquences.
  (1) Il est \emph{insensible au flot} : une écriture de champ faite sur une
  branche est lue sur l'autre. (2) \code{insert} écrase : la ré-exécution de
  \code{const o = {...}} remplace l'objet et efface les écritures de champ
  accumulées sur ce site ; une fermeture ré-allouée remplace ses captures.
  (3) Les sites d'une variable ne sont jamais retirés :
  \code{let f = () => a(); f = () => b();} fait exécuter les deux corps à
  \code{f()}, ce qui sur-approxime. Enfin, la doc de \code{Heap}
  (\og Populated by \code{eval_expr} for \code{FnLit}/\code{ObjectLit}/\code{ArrayLit}
  nodes\fg{}) est inexacte : c'est l'interpréteur qui peuple le tas, et aucun
  tableau n'y entre jamais.
\end{piege}

% ===========================================================================
\section{La valeur d'un mémo}
\label{sec:transfert-stores-memo}

\code{useMemo(f, deps)} et \code{useCallback(f, deps)} rendent la même
référence tant qu'aucune dep ne change au sens d'\code{Object.is}. Leur
valeur abstraite n'est donc pas celle du corps \code{f}, qui n'est jamais
évalué : c'est une \emph{référence} dont la stabilité est le join des
stabilités des deps. Le \cref{chap:statevalue-stabilite} a cité et commenté
\code{recompute_memo} (\src{src/domains/transfer/state_value.rs}{56}{101}) ;
on en retient la table de décision, vérifiée sur
\file{/tmp/ch11/ex17.tsx} (\cref{tab:transfert-stores-memo}).

\begin{table}[htbp]\centering\small
\begin{tabular}{p{48mm}p{24mm}p{36mm}p{26mm}}
\toprule
mémo & deps & valeur (sonde) & raison \\
\midrule
\code{useMemo(() => ({ k: 1 }), [])} & \code{[]} exact & \code{ref(Stable)} & épinglé au montage \\
\code{useMemo(() => ({ k: 1 }))} & absentes & \code{ref(Unknown)} & rien n'est borné \\
\code{useMemo(() => n * 2, [n])} & \code{n} $\in [0,1]$ & \code{ref(PerRender)} & projection \code{to_stability} \\
\code{useMemo(() => obj.a, [obj])} & état objet & \code{ref(Versioned({1}))} & lecture convertie \\
\code{useMemo(() => p + 1, [p])} & prop & \code{ref(Unknown)} & \code{p} vaut $\Top$ \\
\code{useMemo(() => 1, [{}])} & littéral & \code{ref(PerRender)} & allocation \\
\code{useMemo(() => byObj, [byObj])} & un mémo & \code{ref(Versioned({1}))} & étiquettes propagées \\
\code{useCallback(() => setN(1), [])} & \code{[]} & \code{ref(Stable)} & même calcul \\
\bottomrule
\end{tabular}
\caption{La valeur de huit mémos de \file{/tmp/ch11/ex17.tsx}. Une dep qui
est \emph{syntaxiquement} \code{StateVal(l)} donnerait \code{Versioned}
quelle que soit sa sorte ; dans le code réel, \code{[n]} est un
\code{Var("n")}, et un intervalle non ponctuel se projette en
\code{PerRender} : c'est le faux positif de projection relevé au
\cref{chap:statevalue-stabilite}.}
\label{tab:transfert-stores-memo}
\end{table}

Deux décisions encadrent ce calcul. Des deps illisibles (absentes ou
opaques) donnent \code{ref(Unknown)} et jamais \code{Stable} : le trait
\code{Transfer} l'exige (\og it must not share an answer with a written
\code{[]}, which pins the memo forever\fg{}). Et depuis l'ADR-020
(\og \code{recompute_memo} gets the real stores\fg{}), les deps sont
évaluées contre les vrais stores du point fixe : un mémo qui dépend d'un
mémo lit sa valeur courante, et non un $\Bot$ fabriqué.

% ===========================================================================
\section{L'inter-composants, vu du transfert}
\label{sec:transfert-stores-inter}

Le \cref{chap:inter-composants} décrira l'analyse inter-composants du point
de vue du moteur (racines, registre, cache, graphe d'appels). Vue du
transfert, elle se réduit à deux flux : les props descendent vers l'enfant
par l'évaluation d'un \code{CompApp}, les écritures remontent vers le
parent par le store partagé.

\subsection{Un parent, deux enfants}

\begin{lstlisting}[language=TypeScript,caption={Un enfant connu, un enfant inconnu (\file{/tmp/ch11/ex13.tsx})},label={lst:transfert-stores-parent}]
import { useState, useEffect } from "react";
import { Sheet } from "./ui";

function Child({ onChange }: { onChange: (n: number) => void }) {
  useEffect(() => {
    onChange(42);
  }, []);
  return <span />;
}

export function Parent() {
  const [count, setCount] = useState(0);
  const [open, setOpen] = useState(false);
  return (
    <div>
      <Child onChange={setCount} />
      <Sheet onOpenChange={setOpen} />
      {count}{String(open)}
    </div>
  );
}
\end{lstlisting}

\begin{sortie}
-- ComponentId(0) (iterations=0)
   env[onChange] = Val(setter(#1#0))
-- ComponentId(1) (iterations=1)
   state[0] = number[0, 42]   widened=false   shared=number[42, 42]
   state[1] = ⊤   widened=false   shared=⊥
stats: cache_hits=1 cache_misses=1 depth_capped={} unknown_refs={(ComponentId(1), "Sheet")}
\end{sortie}

\code{ComponentId(0)} est \code{Child}, \code{ComponentId(1)} est
\code{Parent}. La chaîne causale :
\begin{enumerate}
  \item la passe de rendu de \code{Parent} atteint le \code{Return} ;
    \code{exec_expr_effects} lance la pré-passe, qui rencontre
    \code{<Child .../>} et l'évalue : \code{eval_comp_app} ;
  \item la prop \code{onChange} vaut \code{setter(#1#0)}, soit
    \code{SetterVal::One(Parent, 0)} ;
  \item \code{Child} est analysé jusqu'à son propre point fixe ; son effet
    appelle \code{onChange(42)}, dont la valeur est un setter d'un autre
    composant : le second bras de \code{exec_setter_call} écrit
    \code{shared_state[(Parent, 0)]} $\join$= $[42,42]$ ;
  \item au test de convergence, \code{Parent} importe sa tranche : le nouvel
    état n'est pas $\lleq$ l'ancien, on itère ; à la seconde passe, l'enfant
    est trouvé dans le cache (\code{cache_hits=1}) ; \code{count} converge
    sur $[0,42]$ ;
  \item \code{Sheet} n'est pas dans le registre : son setter est
    \emph{havoqué}, \code{open} reçoit $\Top$, et une Info le déclare.
\end{enumerate}

\begin{sortie}
$ reactant check ex13.tsx --info
  Parent  (2 hooks)  ex13.tsx
    info   analysis-limit  component `Sheet` was not found in the analysis registry. Pass its file on the command line to analyse it (FN possible)
    suspended  analysis-limit  4 passing check(s) withheld: the analysis was truncated in this component, so they are not guaranteed
\end{sortie}

\subsection{\code{eval_comp_app}, pas à pas}

\begin{enumerate}
  \item \textbf{Sans contexte inter} (analyse d'un composant seul) : tout
    enfant est inanalysable ; on havoque les setters de ses props et on rend
    \code{ref(Stable)} (\src{src/domains/transfer/state_value.rs}{477}{483}).
  \item \textbf{Résolution} du nom JSX par le registre, en un
    \code{ComponentId} (ADR-040). Un enfant inconnu ou ambigu est compté dans
    les statistiques (qui deviendront des Info \regle{analysis-limit}), ses
    setters sont havoqués, et on rend \code{ref(Stable)}
    (\src{src/domains/transfer/state_value.rs}{489}{518}).
  \item \textbf{Récursion} : si l'enfant est déjà sur la pile d'appel, on
    compte une coupure et on rend \code{ref(Stable)}, \emph{sans} havoc.
  \item \textbf{Props} : \code{eval_props_map} rend une carte
    prop $\mapsto$ \code{EnvVal} ; une prop littéral de fonction est allouée
    dans le tas du parent pour que l'enfant puisse l'exécuter. Le code est
    en \src{src/domains/transfer/state_value.rs}{653}{689}.
  \item \textbf{Cache} : les props aplaties en \code{StateValue} servent de
    clé ; un succès enregistre l'arête du graphe d'appels et rend
    \code{ref(Stable)}, sans ré-analyser l'enfant.
  \item \textbf{Contexte de l'enfant}, puis \textbf{analyse} :
\end{enumerate}

\rustsnippet[label={lst:transfert-stores-child}]{src/domains/transfer/state_value.rs}{559}{580}{le tas et l'environnement initiaux d'un enfant}

L'enfant part d'un environnement vide où son paramètre de props est lié à
un site synthétique frais ; son tas initial contient un objet sous ce site
(la carte des props) et une copie des valeurs de tas des props qui ont des
sites (fermetures, objets). \code{analyze_child} est un pointeur vers
\code{analyze_component_inter} : une analyse complète, avec son propre
point fixe, sous un \code{InterCtx} fils dont la pile d'appel contient le
parent. Le résultat est rangé dans la table \code{results} et dans le
cache, l'arête dans le graphe d'appels, et \code{eval_comp_app} rend, dans
tous les cas, \code{ref(Stable)}.

\begin{exemple}{Une boucle croisée}{transfert-stores-croisee}
  Dans \file{/tmp/ch11/ex15.tsx}, \code{Solo} passe
  \code{onTick={() => setB(b + 1)}} à un \code{Ticker} dont l'effet sans deps
  appelle \code{onTick()}. La fermeture est allouée par \code{eval_props_map}
  (captures : \code{setB}, \code{b}) et copiée dans le tas de l'enfant ; dans
  \code{Ticker}, \code{onTick} a un site, et \code{onTick()} est un B6 : le
  corps s'exécute avec ses captures, et \code{setB(b + 1)} tire le second
  bras de \code{exec_setter_call}, puisque le setter a pour propriétaire
  \code{Solo}. Le parent importe l'écriture, itère, et le CLI conclut :
\begin{sortie}
  Ticker  (1 hooks)  ex15.tsx
    warn   cross-component-infinite-loop  [hook:0]  (line 4:2)  this effect calls `onTick`, a state setter of parent `Solo`. Parent re-renders → child re-renders → effect fires again: infinite loop
\end{sortie}
\end{exemple}

\subsection{Le havoc des setters qui s'échappent}

Un setter remis à un enfant \emph{inanalysable} peut être appelé par lui
avec n'importe quel argument, n'importe quand. Ne rien faire
sous-approximerait l'état et fabriquerait des conclusions \og l'état est
stable\fg{}. D'où :

\rustsnippet[label={lst:transfert-stores-havoc}]{src/domains/transfer/state_value.rs}{265}{297}{joindre $\Top$ dans les slots dont le setter s'échappe}

Un setter nu passé en prop est trouvé par sa valeur, qui porte son
propriétaire. Les setters cachés (dans un spread, dans une fermeture qui
appelle un setter, dans une fermeture du tas) sont cherchés par
\code{collect_escaping_setters}
(\src{src/domains/transfer/state_value.rs}{299}{367}), qui parcourt les corps
de fonctions atteignables depuis la valeur de chaque prop. Sa terminaison
est garantie par \emph{identité}, pas par un budget :

\rustsnippet{src/domains/transfer/state_value.rs}{385}{392}{un corps n'est parcouru qu'une fois par environnement capturé}

La clé est la paire (adresse du CFG, adresse des captures) : le même corps
atteint sous d'autres captures résout ses callees autrement. Le commentaire
du code garde la trace de l'ancienne version, un \code{depth > 4} qui était
\og a budget doing a cycle guard's job\fg{} : un setter caché derrière une
cinquième fermeture était silencieusement manqué. Le parcours a néanmoins
une limite, qui est le défaut D13.

\begin{decision}[ADR-012 --- inlining descendant et store partagé]
  L'ADR-012 choisit l'\emph{inlining descendant} : le parent analyse l'enfant
  au point d'application, avec les props abstraites, plutôt que de calculer
  des résumés paramétriques ascendants, jugés \og without justified gain at
  the current stage\fg{}. La sensibilité au contexte passe par l'égalité
  abstraite des props (le cache). Le flux est bidirectionnel : props vers le
  bas, setters vers le haut par un \code{SharedStateStore} passé à toutes les
  analyses, sans couche de point fixe au niveau du programme. Écarts avec le
  code : \code{ComponentSetter} est devenu la composante \code{setter} du
  produit (ADR-015), les clés sont des \code{ComponentId} internés
  (ADR-040), la récursion rend \code{ref(Stable)} sans havoc. Limites
  acceptées : imports hors registre ($\Top$ et Info), récursion (Info),
  composants dynamiques \code{const C = cond ? A : B} non analysés
  (\issue{63}, fermée \emph{wontfix}).
\end{decision}

\begin{piege}
  Trois raccourcis de \code{eval_comp_app} ont des conséquences.
  (1) La valeur d'un élément est toujours \code{ref(Stable)}, y compris pour
  un enfant inconnu : elle n'encode pas l'échec. (2) La récursion ne
  havoque rien : un composant récursif qui passe à son instance intérieure
  un \emph{autre} setter que celui qu'il a reçu sous-approxime ; l'Info de
  récursion le déclare (\og recursive component reference \code{Tree} is not
  followed, so its props are treated as unknown\fg{}), mais son texte ne dit
  pas ce que fait le code, qui n'évalue ni ne havoque aucune prop.
  (3) La clé du cache et la table des résultats oublient une partie du
  contexte (défauts D4 et D7).
\end{piege}

% ===========================================================================
\section{La soundness en pratique : garanties et défauts}
\label{sec:transfert-stores-soundness}

\subsection{Les garanties par construction}

La soundness (\cref{def:soundness}) exige que chaque fonction de transfert
sur-approxime la sémantique concrète. Le code la garantit par une poignée de
mécanismes, que l'on peut maintenant récapituler
(\cref{tab:transfert-stores-garanties}).

\begin{table}[htbp]\centering\small
\begin{tabular}{p{62mm}p{42mm}p{40mm}}
\toprule
mécanisme & garantie & où \\
\midrule
variable inconnue $\Rightarrow$ $\Top$ ; clé d'un seul côté du join $\Rightarrow$ $\Top$ & valeurs & \code{AbstractEnv::lookup}, \code{join_stabs} \\
setter $\Rightarrow$ join dans le slot & écritures may & \code{StateStore::update} \\
setter vers un enfant inconnu $\Rightarrow$ $\Top$ dans le slot & pas de \og stable\fg{} fabriqué & \code{havoc_setter_props} \\
setters échappés cherchés par identité, pas par budget & pas de setter manqué par profondeur & \code{collect_escaping_setters} \\
appel $\Rightarrow$ $\Top$, sauf allocateurs certains & valeurs & \code{returns_fresh_reference} \\
\code{/}, et \texttt{\%} ou \code{**} à \code{NaN} possible $\Rightarrow$ $\Top$ & pas de \code{NaN} dans un intervalle & \code{eval_binop} \\
opérande $\Bot$ $\Rightarrow$ $\Bot$ & un chemin mort reste mort & \code{eval_bitwise}, \code{as_arith} \\
paramètres de callbacks $\Rightarrow$ $\Top$ ; retour d'un corps bouclant $\Rightarrow$ $\Top$ & valeurs & \code{exec_callbacks_depth}, \code{exec_body_impl} \\
écriture de champ $\Rightarrow$ join & mutations may & \code{MemberWrite} \\
deps illisibles $\Rightarrow$ \code{ref(Unknown)}, jamais \code{Stable} & mémos & \code{recompute_memo} \\
\bottomrule
\end{tabular}
\caption{Où la soundness est garantie par construction.}
\label{tab:transfert-stores-garanties}
\end{table}

À côté, des faux négatifs sont \emph{acceptés et déclarés} : un callee
\code{Unknown} non descendu (ADR-009, \issue{19}), le plafond de profondeur
(Info), un enfant inconnu (Info), la récursion (Info), la valeur portée par
une boucle dans un callback (\issue{21}). Le premier est le seul qu'aucune
Info ne signale.

Reste ce que l'invariant du projet interdit : les faux négatifs
\emph{non déclarés}. Les quatorze défauts qui suivent ont été reproduits au
commit \snapshotCommit{} en rédigeant les notes de ce chapitre ; aucun
n'était décrit dans le tracker à cette date (l'issue fermée \issue{130}
traitait l'équivalent de D2 pour la relation \relation{slot_writers}, pas
pour le store). Ce sont des constats, pas des diagnostics confirmés par le
mainteneur. Chacun est une vraie boucle infinie ou une vraie dep manquante
en React, que l'analyseur déclare sûre. Pour chacun, l'exercice est
d'identifier l'hypothèse violée : c'est toujours l'une des garanties
ci-dessus qui n'est pas tenue jusqu'au bout. Les quatorze cas ont été
rejoués pour cette relecture, au commit \snapshotCommit{}, sur des copies
fraîches dans \file{/tmp/ch11/} : chacun reproduit exactement la valeur ou
le silence du CLI annoncé plus bas. Restent constats, non diagnostics : la
question de savoir si D8 et D12 relèvent d'une décision de conception
tacite plutôt que d'un défaut reste au mainteneur.

\subsection{D1 : un corps qui voit l'état d'un bloc en retard}

\begin{lstlisting}[language=TypeScript,caption={Trois callbacks à deux blocs (\file{/tmp/ch11/ex3.tsx}, extrait de l'effet)},label={lst:transfert-stores-d1}]
    setTimeout(() => {
      const v = 5;
      if (flag) {
        setN(v);
      }
    }, 10);
    setTimeout(() => {
      const w = 7;
      setM(w);
    }, 10);
    setTimeout(() => {
      const s = setK;
      if (flag) {
        s(3);
      }
    }, 10);
\end{lstlisting}

\begin{sortie}
-- OneBlockBehind (intra, iterations=1)
   state[0] = ⊤   widened=false
   state[1] = number[0, 7]   widened=false
   state[2] = number[0, 0]   widened=false
\end{sortie}

Le callback en ligne droite est exact (\code{m} $\in [0,7]$). Les deux
autres ne le sont pas : \code{n} lit \code{v} comme $\Top$, et l'alias
\code{s} est perdu, si bien que l'écriture de \code{k} \emph{disparaît}. La
cause est dans \code{exec_body_impl} (\cref{lst:transfert-stores-body}) :
chaque terminateur écrit \code{env_at[next]}, qui est donc
l'environnement \emph{d'entrée} de \code{next} ; mais l'environnement d'un
bloc non initial est calculé comme le join des \code{env_at[p]} de ses
prédécesseurs \code{p}, c'est-à-dire de leurs environnements d'\emph{entrée}
(\cref{fig:transfert-stores-d1}). On attendrait
\code{env_at.get(&bid)}. Une valeur perdue retombe à $\Top$ (sound) ; un
alias de setter, un site ou une liaison de \code{useCallback} perdus
suppriment une écriture (FN). C'est aussi ce qui met \code{UpdaterGuarded}
de l'\cref{ex:transfert-stores-ternaire} à $\Top$ : le bloc de jonction lit
\code{__t0} dans l'environnement d'entrée de \code{B1} et \code{B2}, où il
n'est pas lié. Les tests unitaires ne le voient pas : leurs corps à
plusieurs blocs ne lisent que des variables de l'environnement initial.

\begin{figure}[htbp]\centering
\begin{tikzpicture}[font=\scriptsize]
  \node[cfgbloc] (b0) at (0,0) {B0\\ let s = setK\\ => branch flag ? B1 : B2};
  \node[cfgbloc] (b1) at (-2.2,-2) {B1\\ expr s(3)\\ => jump B3};
  \node[cfgbloc] (b2) at (2.2,-2) {B2\\ => jump B3};
  \node[cfgbloc] (b3) at (0,-4) {B3\\ => return undefined};
  \draw[fleche] (b0) -- (b1); \draw[fleche] (b0) -- (b2);
  \draw[fleche] (b1) -- (b3); \draw[fleche] (b2) -- (b3);
  \node[etiquette,align=left,text width=46mm,anchor=west] at (2.3,0) {entrée de B0 : $E$, où \code{s} n'est pas lié ;\\ sortie de B0 : $E$ plus \code{s}, setter de \code{k}};
  \node[etiquette,align=left,text width=40mm,anchor=east,text=ra-amber] at (-3.4,-2) {B1 lit \code{env_at[B0]}, soit $E$ :\\ \code{s} est inconnu, \code{s(3)} n'écrit rien. Il devrait lire \code{env_at[B1]}, la sortie de B0};
  \node[etiquette,align=left,text width=40mm,anchor=west] at (3.2,-2) {\code{env_at[B1]} et \code{env_at[B2]} valent la sortie de B0 : écrits par le terminateur de B0, jamais lus par B1 ni B2};
\end{tikzpicture}
\caption{Le défaut D1 sur le troisième callback de \cref{lst:transfert-stores-d1}.
\code{env_at} est indexé par le bloc \emph{destinataire}, mais chaque bloc
lit la table sous les noms de ses \emph{prédécesseurs} : il voit l'entrée de
son prédécesseur, soit l'état \og d'un bloc en retard\fg{}.}
\label{fig:transfert-stores-d1}
\end{figure}

\subsection{D2 et D3 : où et comment un setter écrit}

\textbf{D2, le setter hors position d'effet.} L'\cref{inv:setter-positions}
dit où un setter écrit. Trois effets sans deps, identiques au sens de React,
le mettent à l'épreuve (\file{/tmp/ch11/ex6.tsx}) :
\code{setA(a + 1);} (\code{StmtPos}), \code{flag && setA(a + 1);}
(\code{AndPos}), \code{console.log(setA(a + 1));} (\code{ArgPos}).

\begin{sortie}
-- StmtPos (intra, iterations=3)
   state[0] = number[0, inf]   widened=true
-- AndPos (intra, iterations=0)
   state[0] = number[0, 0]   widened=false
-- ArgPos (intra, iterations=0)
   state[0] = number[0, 0]   widened=false
\end{sortie}

Le CLI signale \code{StmtPos} (\regle{infinite-loop}) et déclare les deux
autres propres. Dans \code{AndPos}, le lowering a rangé l'appel dans
\code{__t0 := setA((a Add 1))} (un \code{Assign}) ; dans \code{ArgPos},
l'appel est un argument. Aucun des deux n'atteint \code{exec_setter_call}.
Pire : si l'argument est un \emph{functional updater} en position de
liaison (\code{const r = setB((x) => ...)}), la pré-passe classe le callee
\code{Setter} et ne descend pas le littéral, en comptant sur
\code{exec_setter_call}, qui n'est pas appelé : ni l'écriture ni les effets
du corps ne sont vus. \file{docs/limitations.md} affirme qu'un appel de
setter \og in any expression position [\ldots] is seen\fg{} : c'est vrai pour
la relation \relation{slot_writers}, pas pour le store d'état.

\textbf{D3, l'updater en mode programme.} L'\cref{ex:transfert-stores-updater}
annonçait que le CLI ne retrouve pas le widening de \code{UpdaterNoDeps}.
La sonde, en mode programme, le montre :
\begin{sortie}
-- ComponentId(0) (iterations=1)
   state[0] = ⊤   widened=false   shared=ref(PerRender)
\end{sortie}
Le premier bras de \code{exec_setter_call} exécute l'updater correctement ;
le second, qui tire aussi pour un setter local dès que \code{ctx.inter} est
présent, évalue l'argument \code{FnLit} comme une valeur :
\code{ref(PerRender)}, la fermeture elle-même. Cette fermeture est écrite
dans \code{shared_state[(composant, 0)]}, ré-importée par \code{slice} dans
l'état ; le paramètre \code{x} devient \og nombre ou référence\fg{},
\code{x + 1} vaut $\Top$, et plus aucun widening n'a lieu. Sur la copie du test
\code{functional_updater_without_deps_is_flagged}
(\file{/tmp/ch11/ex8.tsx}, \code{useEffect(() => { setCount(c => c + 1); })}),
le CLI répond ✓ \code{1 file(s) no issues found.} alors que le test
d'intégration, qui passe par le chemin intra-composant, est vert. C'est un
défaut d'\emph{interaction} : chaque bras est raisonnable isolément. Une
correction possible — ne tirer le second bras que si le premier n'a pas
tiré, ou si le propriétaire n'est pas \code{ctx.component} — reste à
trancher par le mainteneur : elle touche la sémantique du second bras pour
tout appel, pas seulement l'updater.

\subsection{D4 et D7 : le contexte d'un enfant, oublié}

Les deux défauts viennent de la même hypothèse : l'analyse d'un enfant ne
dépendrait que des valeurs aplaties de ses props.

\textbf{D4.} Dans \file{/tmp/ch11/ex14.tsx}, \code{Board} rend deux
\code{Ticker} (\cref{ex:transfert-stores-croisee}) :
\code{onTick={() => setA(1)}} puis \code{onTick={() => setB(b + 1)}}. Les
deux fermetures valent \code{ref(PerRender)} : la clé du cache est la même,
le second site d'appel touche le cache, et l'enfant n'est jamais analysé
avec la fermeture qui boucle (\code{cache_hits=15 cache_misses=1} pour tout
le fichier). Le CLI est muet ; le même fichier où le second enfant est un
composant \code{Ticker2} identique mais distinct donne
\regle{cross-component-infinite-loop}. La clé ignore les \emph{corps} des
fermetures, c'est-à-dire le tas.

\textbf{D7.} \code{Counter({ step })} écrit \code{setN(n + step)} dans un
effet sans deps. Rendu seul avec \code{step={1}}, il est signalé
(\regle{infinite-loop}). Rendu deux fois,
\code{<Counter step={1} /><Counter step={0} />} (\file{/tmp/ch11/ex22.tsx}),
le CLI répond \code{Counter (2 hooks)} ✓ : la seconde analyse a écrasé la
première dans la table \code{results}
(\src{src/domains/transfer/state_value.rs}{583}{586}), et les règles ne lisent
que la dernière. La table n'a qu'une entrée par \code{ComponentId}, alors
qu'un composant a autant d'instances que de sites d'appel.

\subsection{D5, D6, D9, D11, D13, D14 : le tas}

Le tas concentre six défauts, tous des cas où un site, un membre ou un
corps n'est pas suivi jusqu'au bout.
\begin{itemize}
  \item \textbf{D5, lecture partielle d'un membre.} Avec
    \code{const o = flag ? { f: 1 } : { ...p }}, \code{o.f} est annoncé
    \code{number[1, 1]} alors qu'il peut valoir \code{p.f} ; même chose
    pour \code{q?.f} avec \code{q} nullable, et pour \code{r.f} quand l'autre
    branche est \code{{ g: 2 }} (\file{/tmp/ch11/ex18.tsx}). Conséquence
    observée, dans \file{/tmp/ch11/ex19.tsx} : un
    \code{useEffect(() => console.log(v), [])} où \code{v} vient d'un tel
    ternaire (\code{PartialDeps}) n'est pas signalé par
    \regle{missing-deps}, alors que la variante sans ternaire l'est
    (\code{SpreadOnlyDeps}, même fichier).
  \item \textbf{D6, objets imbriqués et appels de méthodes.} B6 ne traite
    que les callees \code{Var}, et un littéral imbriqué n'a pas de site :
    \code{o.f()} et \code{o.inner.f} manquent leur setter (vu en
    \cref{sec:transfert-stores-tas}). Avec
    \code{o = { f: () => setN(n + 1) }} dans un effet sans deps,
    \code{g = o.f; g()} mène \code{n} au widening, \code{o.f()} non
    (\file{/tmp/ch11/ex12.tsx}).
  \item \textbf{D9, champ nouveau écrit sur une branche.}
    \code{if (flag) { o.g = 2 } const v = o.g;} donne \code{v = [2, 2]},
    sans \code{undefined} (\file{/tmp/ch11/mw.tsx}) : le cas
    \code{(None, new)} de \code{MemberWrite} n'ajoute pas \code{undefined},
    et le tas insensible au flot fait voir l'écriture au chemin qui ne l'a
    pas faite. Un champ \emph{existant}, lui, est bien joint
    (\code{number[1, 1]|string{"x"}}).
  \item \textbf{D11, fermeture membre d'un objet prop.} Seuls les sites
    \emph{directs} des props sont copiés dans le tas de l'enfant
    (\cref{lst:transfert-stores-child}). Avec
    \code{bag = { run: () => setN(n + 1) }} passé en prop, l'enfant qui
    appelle \code{bag.run} trouve un site sans valeur de tas et n'exécute
    rien ; la même fermeture passée directement (\code{run={run}}) donne
    \regle{cross-component-infinite-loop} (\file{/tmp/ch11/deep.tsx}).
  \item \textbf{D13, deux sauts de fermeture.} Dans un corps parcouru,
    \code{collect_escaping_setters} ne suit pas un appel vers une
    \emph{autre} fermeture : avec
    \code{inner = (v) => setN(v); outer = (v) => inner(v)} passé dans
    \code{<Sheet cfg={{ cb: outer }}/>}, \code{n} reste $[0,0]$, alors qu'avec
    un seul saut il vaut $\Top$ (\file{/tmp/ch11/esc2.tsx}). Ce défaut-là est
    couvert globalement par l'Info \og component \code{Sheet} was not
    found\fg{}, qui suspend les vérifications ; la valeur calculée reste
    fausse.
  \item \textbf{D14, l'exécuteur de \code{new Promise}.} La pré-passe ne
    classe que les \code{Call} ; un \code{New} passe par la descente
    générique, qui n'entre pas dans un littéral de fonction. L'exécuteur de
    \code{new Promise((resolve) => { setN(n + 1); ... })}, que JavaScript
    exécute \emph{synchroniquement}, n'est jamais exécuté : \code{n} reste
    $[0,0]$ et \regle{infinite-loop} se déclare vérifiée
    (\file{/tmp/ch11/np.tsx}). Le cas fréquent
    \code{new Promise(r => setTimeout(r, 100)).then(...)} n'est pas touché.
\end{itemize}

\subsection{D8, D10, D12 : trois modélisations}

\textbf{D8.} \code{const el = <span />} vaut \code{ref(Stable)} ;
\code{useEffect(..., [el])} n'est pas signalé par
\regle{always-unstable-deps}, alors que \code{[obj]} avec un littéral
d'objet l'est (\file{/tmp/ch11/ex24.tsx}). En React, \code{el} est un objet
neuf à chaque rendu. Ce pourrait être une décision implicite (éviter des FP
sur \code{children}) ; aucun ADR ni commentaire ne la justifie.

\textbf{D10.} \code{const s = c ? setA : setB; s(1);} n'écrit que
\code{a} (\file{/tmp/ch11/sj.tsx}) : le join à priorité gauche des liaisons
de setter garde \code{setA}, et la valeur jointe, \code{SetterVal::Top}, ne
satisfait pas \code{as_setter()}, si bien que le second bras ne tire pas
non plus. Le résultat dépend de l'ordre des opérandes du join.

\textbf{D12.} \code{useReducer} ignore le réducteur :
\begin{sortie}
-- Red (intra, iterations=1)
   state[0] = number[0, 1]   widened=false
   env[dispatch] = Val(setter(#4294967295#0))
\end{sortie}
Avec \code{reducer = (s, a) => s + a} et
\code{useEffect(() => { dispatch(1); })} (\file{/tmp/ch11/red.tsx}), l'action
$1$ est écrite telle quelle, le store converge sur $[0,1]$ et le CLI répond
✓. En React, \code{s} vaut $1, 2, 3, \ldots$ et l'effet relance le
rendu indéfiniment. La cause est au lowering (le réducteur est sauté,
\cref{chap:lowering-cfg}), l'effet ici.

\subsection{Synthèse}

\begin{table}[htbp]\centering\small
\begin{tabular}{p{9mm}p{52mm}p{56mm}p{22mm}}
\toprule
& défaut & hypothèse violée & exemple \\
\midrule
D1 & \code{exec_body_impl} lit l'entrée des prédécesseurs & propagation d'un CFG & \file{ex3.tsx} \\
D2 & setter hors \code{ExprStmt}/\code{Return} & \og tout appel qui peut avoir lieu est vu\fg{} & \file{ex6.tsx} \\
D3 & second bras sur un updater local & un bras n'annule pas l'autre & \file{ex8.tsx} \\
D4 & clé de cache aplatie & contexte = props aplaties & \file{ex14.tsx} \\
D5 & membre joint sur les seuls sites qui l'ont & un site sans le membre compte & \file{ex18.tsx} \\
D6 & B6 limité aux \code{Var}, littéraux imbriqués sans site & tout corps atteignable est suivi & \file{ex11.tsx} \\
D7 & \code{results} écrasé par site d'appel & une instance par site & \file{ex22.tsx} \\
D8 & élément JSX \code{Stable} & fraîcheur d'une allocation & \file{ex24.tsx} \\
D9 & champ nouveau sans \code{undefined} & écriture may, lecture ailleurs & \file{mw.tsx} \\
D10 & liaison de setter à priorité gauche & join des liaisons latérales & \file{sj.tsx} \\
D11 & sites des membres de props non copiés & tas de l'enfant complet & \file{deep.tsx} \\
D12 & réducteur ignoré & sémantique de \code{dispatch} & \file{red.tsx} \\
D13 & un seul saut de fermeture & transitivité de l'échappement & \file{esc2.tsx} \\
D14 & \code{new} jamais classé & exécuteur synchrone & \file{np.tsx} \\
\bottomrule
\end{tabular}
\caption{Les quatorze défauts, tous reproduits sur \file{/tmp/ch11/}. D13
est couvert par une Info globale ; les autres sont des faux négatifs non
déclarés.}
\label{tab:transfert-stores-defauts}
\end{table}

\begin{remarque}
  Plusieurs de ces défauts se corrigent par une ligne ou par un déplacement
  de responsabilité, dans l'esprit du principe \og pas de workaround\fg{} :
  lire \code{env_at[bid]} dans \code{exec_body_impl} (D1) ; tirer les
  setters depuis un point central qui voit toute position d'expression,
  sans double tir (D2) ; faire entrer les sites des props dans la clé du
  cache (D4). Ils sont proposés en exercice.
\end{remarque}

% ===========================================================================
\begin{resume}
\begin{itemize}
  \item L'interpréteur abstrait d'un pas tient en deux fonctions. La
    première, \code{eval_state_value}, est une récursion structurelle qui
    rend une \code{StateValue} sans jamais écrire l'état ; elle vit dans
    \file{transfer/state_value.rs}. La seconde, \code{exec_full_stmt}, dans
    \file{interp/interpreter.rs}, exécute une instruction en deux temps :
    pré-passe des callbacks, puis cœur. Le moteur les atteint par le trait
    \code{Transfer}, dont l'unique implémentation est
    \code{StateValueTransfer}.
  \item L'environnement \code{AbstractEnv} a quatre tables. Une variable
    inconnue vaut $\Top$ ; une liaison de setter, de callback ou de site
    inconnue ne vaut \emph{rien}, et sa perte coûte une écriture.
  \item Un appel vaut $\Top$ sauf allocateur certain ; un littéral vaut
    \code{ref(PerRender)} ; un élément JSX \code{ref(Stable)}. Le lowering a
    déjà transformé ternaires, \code{&&}, \code{||}, \code{??} en diamants et
    les gabarits en additions.
  \item Quatre stores (\cref{def:store}) : \code{StateStore} ($\Bot$, join,
    sujet du point fixe, amorcé par les initialiseurs), \code{MemoStore}
    ($\Top$, écrasement, recalculé après chaque rendu),
    \code{SharedStateStore} ($\Bot$, join, un par programme, importé par
    \code{slice}), \code{Heap} (par site, écrasement, insensible au flot).
  \item Un setter écrit par mise à jour faible, en trois positions
    seulement ; un \emph{functional updater} lit le join du store. La lecture
    d'un état est convertie en \code{Versioned} en un seul point.
  \item Les callbacks sont classés par ce qui les déclenche
    (\code{Setter}, \code{InCycle}, \code{Subscription}, \code{Unknown}) ;
    seuls les \code{InCycle} sont exécutés, par un second interpréteur de CFG
    en une passe, jusqu'à la profondeur 3. \code{Unknown} non descendu est un
    FN accepté (ADR-009).
  \item Le tas relie un nom à un corps (B5, B6) et un objet à ses membres
    (\issue{88}) ; l'inlining d'un enfant copie les props dans un tas neuf, et
    les écritures remontent par le store partagé (ADR-012).
  \item Quatorze défauts (D1 à D14) montrent où ces garanties s'arrêtent :
    propagation des corps, positions des setters, contexte des enfants,
    membres et sites non suivis.
\end{itemize}
Le \cref{chap:cfg-analyzer-dominance} décrira l'autre interpréteur de CFG,
\code{analyze_cfg}, avec sa worklist, ses raffinements et ses arcs arrière ;
le \cref{chap:point-fixe-composant} la boucle qui enchaîne rendu, mémos,
effets et handlers jusqu'au point fixe.
\end{resume}

% ===========================================================================
\section*{Exercices}
\addcontentsline{toc}{section}{Exercices}
\label{sec:transfert-stores-exercices}

\begin{exercice}{Calculer à la main}{transfert-stores-calculs}
  Donner la valeur abstraite de : (a) \code{a + 1} avec
  $a = \code{number[0, 3]|null}$ ; (b) \code{x & 7} avec $x = \Top$ ;
  (c) \code{typeof null} ; (d) \code{+s} avec $s = \code{string{"5", "6"}}$ ;
  (e) \code{!n} avec $n = \code{number[0, 0]}$ ; (f) \code{"a" + 1} ;
  (g) \code{+""}.
\end{exercice}
\begin{solution}
  (a) \code{as_arith} accepte \code{null}, qui vaut $0$ : l'opérande devient
  $[0,3] \join [0,0] = [0,3]$, et la somme $[1,4]$. (b) Le masque est une
  constante positive : $[0,7]$, quel que soit $x$. (c) Une seule sorte
  habitée, \code{null} : \code{string{"object"}}. (d) Les deux chaînes se
  lisent comme des nombres finis : \code{number[5, 6]}. (e) L'opérande n'est
  pas un booléen pur : $\Top$ (et non \code{boolean}). (f) Chaîne plus
  nombre : $\Top$. (g) \code{"".trim().parse::<f64>()} échoue en Rust :
  $\Top$, alors que JavaScript donne $0$ ; c'est sound.
\end{solution}

\begin{exercice}{Un gabarit}{transfert-stores-gabarit}
  Pour \code{const t = `id-${n}`} avec \code{n} un état numérique, donner
  l'IR produite et la valeur de \code{t}. Proposer une amélioration sound.
\end{exercice}
\begin{solution}
  La sonde d'impression donne \code{let t = (("id-" Add n) Add "")}, et la
  sonde de valeurs donne $\Top$ pour \code{t} : \code{"id-" + n} tombe dans le
  cas \og tout le reste\fg{} de \code{+}. Or en JavaScript, \code{+} dont
  un opérande est une chaîne rend toujours une chaîne : rendre
  \code{str_top()} dès qu'un côté est une chaîne pure serait sound et plus
  précis (la composante nombre resterait $\Bot$).
\end{solution}

\begin{exercice}{Dérouler un updater}{transfert-stores-updater-main}
  Le slot $0$ vaut $[0,2]$. On exécute \code{setN(c => c + 1)}. Quelle
  valeur est écrite, que vaut le slot ensuite, et pourquoi contient-il encore
  $0$ ?
\end{exercice}
\begin{solution}
  \code{c} est lié à \code{ctx.state.get(0)}, soit $[0,2]$ ; le corps rend
  $[1,3]$, qui est la valeur écrite ; le slot devient
  $[0,2] \join [1,3] = [0,3]$. Il contient $0$ parce que la mise à jour est
  faible : le store range le join de \emph{toutes} les valeurs que l'état a
  pu prendre, initialiseur compris, pas sa valeur courante.
\end{solution}

\begin{exercice}{Classer des callees}{transfert-stores-classer}
  Donner la \code{TriggerClass} de \code{promise.then(f)},
  \code{arr.forEach(f)}, \code{el.addEventListener("x", f)},
  \code{myLib.subscribe(f)}, \code{Array.from(xs, f)},
  \code{router.from(f)}, et dire si \code{f} est exécuté.
\end{exercice}
\begin{solution}
  \code{then} : \code{InCycle}, exécuté. \code{forEach} : \code{InCycle},
  exécuté. \code{addEventListener} : \code{Subscription}, non exécuté ici,
  mais relevé comme handler et analysé comme point d'entrée s'il est dans un
  effet. \code{subscribe} : \code{Unknown}, non exécuté (FN accepté,
  \issue{19}). \code{Array.from} : \code{InCycle}, le receveur étant
  \code{Array}. \code{router.from} : \code{Unknown}, car \code{from} n'est
  reconnu que sur le receveur \code{Array}.
\end{solution}

\begin{exercice}{Le plafond de profondeur}{transfert-stores-profondeur}
  Montrer que, dans \cref{lst:transfert-stores-callbacks}, \code{f4()}
  n'atteint pas \code{setD}. Où l'utilisateur en est-il averti ? Pourquoi
  ne le serait-il pas si le composant était analysé seul ?
\end{exercice}
\begin{solution}
  L'effet s'exécute à la profondeur $0$ ; \code{f4()} est un B6 qui exécute
  le corps de \code{f4} à la profondeur $1$, celui de \code{f3} à $2$, celui
  de \code{f2} à $3$. Dans ce dernier, la pré-passe de l'instruction
  \code{f1()} trouve $d \geq 3$ et s'arrête : \code{f1} n'est jamais
  exécuté. En mode programme, l'arrêt est enregistré dans
  \code{callback_depth_capped}, d'où l'Info \regle{analysis-limit} et la
  suspension des vérifications. L'enregistrement est gardé par
  \code{if let Some(inter) = ctx.inter} : une analyse intra-composant (tests,
  phase 2 d'\code{analyze_program}) subit le plafond en silence.
\end{solution}

\begin{exercice}{L'environnement vide}{transfert-stores-bottom}
  Calculer \code{AbstractEnv::bottom().join(&e)} pour
  $e = \{x \mapsto [1,1]\}$. Pourquoi \code{analyze_cfg} n'est-il pas
  affecté ?
\end{exercice}
\begin{solution}
  $x$ n'est présent que du côté droit : il vaut $\Top$. Le résultat est
  $\{x \mapsto \Top\}$, pas $e$ : l'environnement vide n'est pas neutre.
  \code{analyze_cfg} n'y est pas exposé parce qu'il insère le premier
  environnement reçu par un bloc sans join (\code{None => outgoing_env}),
  et ne joint qu'ensuite des environnements effectivement calculés.
\end{solution}

\begin{exercice}{Corriger D1}{transfert-stores-d1}
  Localiser dans \cref{lst:transfert-stores-body} la ligne responsable de D1,
  écrire la correction, et décrire le test unitaire qui l'empêcherait de
  revenir.
\end{exercice}
\begin{solution}
  Pour un bloc non initial, l'environnement est le join des
  \code{env_at.get(p)} des prédécesseurs
  (\src{src/domains/interp/interpreter.rs}{419}{425}). Or les terminateurs
  écrivent déjà dans \code{env_at[next]} le join des sorties des
  prédécesseurs de \code{next} : il suffit de lire
  \code{env_at.get(&bid).cloned().unwrap_or_else(AbstractEnv::bottom)} pour
  tout bloc. Le test : un corps \code{() => { const s = setK; if (c) { s(3); } }}
  exécuté par \code{exec_body} doit écrire $[3,3]$ dans le slot de
  \code{setK} ; aujourd'hui, il n'écrit rien. Un second test sur un
  updater ternaire (\cref{ex:transfert-stores-ternaire}) doit rendre un
  intervalle et non $\Top$.
\end{solution}

\begin{exercice}{Une correction centrale de D2}{transfert-stores-d2}
  Dans l'esprit du principe \og pas de workaround\fg{} de \file{CLAUDE.md},
  proposer un endroit unique où tirer les appels de setter pour qu'ils
  soient vus dans toute position d'expression, sans double tir.
\end{exercice}
\begin{solution}
  Une réponse possible : la pré-passe visite déjà \emph{tous} les appels, en
  toute position, et classe \code{Setter} les callees liés à un setter. C'est
  là qu'il faudrait appeler \code{exec_setter_call}, et le retirer de
  \code{exec_expr_effects} et du \code{Return} d'\code{exec_body_impl}, pour
  qu'un appel ne soit tiré qu'une fois. Le bras \code{Setter} y exécuterait
  aussi l'updater, ce qui corrige la variante \og updater en position de
  liaison\fg{}. Il faut ensuite s'assurer qu'aucun appel n'est visité deux
  fois par la pré-passe au sein d'une même instruction, et que les rejeux
  post-convergence du moteur restent cohérents. Ce n'est qu'une piste, non
  implémentée ni testée ; le mainteneur tranchera.
\end{solution}

\begin{exercice}{Une clé de cache qui voit D4}{transfert-stores-d4}
  Sur \file{/tmp/ch11/ex14.tsx}, expliquer pourquoi la boucle n'est pas vue,
  et proposer une clé de cache qui la verrait.
\end{exercice}
\begin{solution}
  Les deux props \code{onTick} valent \code{ref(PerRender)} une fois
  aplaties ; la clé est la même, le second site d'appel touche le cache.
  La carte \code{abstract_props_full} contient pourtant les \emph{sites} des
  fermetures (\code{EnvVal::Loc}). Une clé qui inclurait, pour chaque prop,
  l'ensemble de ses sites distinguerait les deux appels. Il faudrait alors
  que \code{results} soit lui aussi indexé par site d'appel, sans quoi D7
  masquerait la seconde analyse.
\end{solution}
